\chapter{Non-Archimedean psh functions}
\index{plurisubharmonic function!BFJ}
\index{skeletal characterization}
\index{singular semipositive metric}
\label{chap:bfj-psh}

Model functions are finite and piecewise affine.  Singular
plurisubharmonic functions arise only after completing that model class on
every skeleton and imposing compatibility with all retractions.  The two
requirements play different roles: the skeletal condition remembers
positivity, while the retraction inequality determines the values away from
quasimonomial points.

We retain the standing BFJ hypotheses of
\Cref{chap:bfj-model-forms}.  Fix

\begin{equation}
\label{eq:ch21-standing-theta}
            \theta\in\mathcal Z^{1,1}(X),
            \qquad \{\theta\}\in N^1(X)\text{ ample}.
\end{equation}

The form \(\theta\) need not itself be semipositive.  By
\Cref{prop:ch20-positive-representative}, however, it becomes positive after
adding \(\ddc\) of a model function, so the psh class defined below is
nonempty.

\section{Orientation}
\label{sec:ch21-orientation}

\subsection{What a definition of singular potentials must accomplish}
\label{sec:ch21-what-must-be-done}

\Cref{chap:bfj-model-forms} produced a class of potentials that is far too
small for any variational argument: the model functions.  They are finite
everywhere, continuous, and piecewise affine on every dual complex; the
\(\theta\)-psh ones among them form a convex set stable under maxima
(\Cref{prop:ch20-model-max}).  What they are not is \emph{complete}.  Every
theorem obtained by minimizing a functional, by taking a supremum over a
family, or by letting a sequence of potentials degenerate produces a limit
object outside the model class, and no amount of care in the choice of
minimizing family can avoid this.  The situation is exactly that of the
complex theory, where the analogue of a model function is a smooth
potential and the analogue of the completion carried out here is the class
of \(\omega\)-psh functions.

Three requirements shape the definition.

\begin{enumerate}
\item \emph{Positivity must be remembered.}  Whatever the new class is, its
members must still satisfy the numerical positivity that
\Cref{chap:bfj-model-forms} attached to a relatively nef vertical divisor;
otherwise the intersection-theoretic identities on which the whole part
rests will not survive the passage to the limit.
\item \emph{The class must be closed under the limits that actually
occur.}  Decreasing limits and upper envelopes must stay inside.  This
forces the values \(-\infty\) to be admitted, and it forces upper
semicontinuity, since an upper envelope is not upper semicontinuous without
regularization (\Cref{prop:usc-majorant}).
\item \emph{A function must be determined by computable data.}  A psh
function should be reconstructible from its restrictions to dual complexes,
which are finite polyhedra on which everything is a piecewise affine
computation, and ultimately from its values at divisorial points, which are
integers divided by multiplicities.
\end{enumerate}

The definition of \Cref{sec:ch21-definition} meets all three, and each of
its two clauses is responsible for a different one.  Clause~(2), the
skeletal approximation clause, carries requirement~(1): a uniform limit of
restrictions of \(\theta\)-psh model functions inherits every convexity and
slope inequality that positivity imposes on a dual complex.  Clause~(1), the
retraction inequality, carries requirement~(3): together with upper
semicontinuity it forces the reconstruction formula
\eqref{eq:ch21-skeletal-reconstruction}, so that the function is the
decreasing limit of its own skeletal shadows.  Requirement~(2) is then a
theorem, proved in \Cref{sec:ch21-stability}.

\subsection{Standing hypotheses and where each of them is used}
\label{sec:ch21-hypotheses}

Everything in this chapter is proved under the original BFJ hypotheses,
which we restate because their individual roles differ and because later
chapters remove some of them:

\begin{equation}
\label{eq:ch21-hypotheses}
\left\{
\begin{aligned}
 &K\ \text{complete discretely valued},\ v(\pi)=1,\
   \operatorname{char}\widetilde K=0;\\
 &X\ \text{smooth, connected, projective over }K,\ \dim X=n;\\
 &\text{models normal, flat, projective over }R=K^\circ;\\
 &\theta\in\mathcal Z^{1,1}(X)\ \text{with}\ \{\theta\}\ \text{ample}.
\end{aligned}
\right.
\end{equation}

The following table records, for each hypothesis, the single place in this
chapter where it is genuinely used, and what fails without it.  A reader who
wants to know whether a statement survives a weakening should consult this
table first and \Cref{sec:ch21-dictionary} second.

\begin{center}
\small
\begin{tabular}{p{.24\textwidth}|p{.34\textwidth}|p{.34\textwidth}}
\toprule
\textbf{Hypothesis}&\textbf{Where it is used here}
 &\textbf{What fails without it}\\
\midrule
\(K\) discretely valued
 &throughout: the definition quantifies over models over \(R\), and
  \(\varphi_{\mathcal X_s}=1\) fixes the normalization \eqref{eq:ch20-sign}
 &there is no theory of models over a non-discrete valuation ring in the
  form used in \Cref{part:bfj}; the formal-model substitute changes the
  meaning of ``vertical divisor''\\
\(\operatorname{char}\widetilde K=0\)
 &\Cref{thm:ch21-model-consistency}, through the multiplier-ideal package
  of \Cref{sec:ch21-multiplier-input}; and, through
  \Cref{prop:snc-cofinal}, in the cofinality of SNC models
 &the closedness statement \eqref{eq:ch21-model-closedness} is not
  available by this proof; resolution of singularities is not available for
  the cofinality of SNC models\\
\(X\) smooth
 &\Cref{prop:snc-cofinal} and the SNC normal form of
  \Cref{lem:snc-normal-form}, hence the entire skeletal calculus
 &dual complexes need not be simplicial and the monomial description of
  \Cref{lem:monomial-valuation} fails at a singular point of \(X\)\\
\(X\) projective
 &\Cref{thm:ch20-ddc-exact} and every use of a relatively ample class
 &\(N^1(\mathcal X/R)\) loses the ample cone used in
  \Cref{prop:ch20-model-max}\\
\(\{\theta\}\) ample
 &\Cref{prop:ch20-positive-representative}, hence nonemptiness of
  \(\PSH(X,\theta)\) and \Cref{prop:ch21-constants}
 &the class can be empty; \eqref{eq:ch21-reference-class-example} has no
  positive representative to change to\\
\bottomrule
\end{tabular}
\end{center}

Two of these entries deserve emphasis.  Ampleness of \(\{\theta\}\) is a
hypothesis on a \emph{numerical class}, not on the form: \(\theta\) itself
is never assumed semipositive, and \Cref{ex:ch21-reference-class} shows that
assuming it would be a genuine loss of generality rather than a
normalization.  Residue characteristic zero enters at exactly one place in
this chapter, and it is worth naming that place now:
\Cref{thm:ch21-model-consistency}, whose closedness half uses asymptotic
multiplier ideals.  Nothing else in the chapter---not the definition, not
the reconstruction theorem, not the stability calculus, not the polar-set
theory---uses it.

\subsection{Global approximation versus local regularity}
\label{sec:ch21-global-vs-local}

\index{plurisubharmonic function!BFJ versus CLD}
This book contains two constructions that are both called
``plurisubharmonic'', and the reader must keep them apart from the outset.
They are not two presentations of one notion.  They are built from
different data, they are stable under different operations, and the book
proves their agreement exactly on the loci where
\Cref{thm:ch19-cld-bfj-comparison} proves it and nowhere else.

\begin{center}
\small
\begin{tabular}{p{.22\textwidth}|p{.36\textwidth}|p{.34\textwidth}}
\toprule
 &\textbf{BFJ \(\theta\)-psh (this chapter)}
 &\textbf{CLD psh (\Cref{chap:cld-poincare-lelong})}\\
\midrule
defining mechanism
 &\emph{global approximation}: uniform approximation on every dual complex
  by restrictions of \(\theta\)-psh model functions
  (\Cref{def:ch21-bfj-psh})
 &\emph{local regularity}: membership in a sheaf generated from smooth
  functions with \(\dprimeddprime u\ge0\) by local decreasing limits
  (\Cref{def:ch15-hierarchy})\\
ambient datum
 &a projective \(K\)-variety \(X\) and a closed model form \(\theta\);
  there is no notion over a proper open subset
 &an arbitrary boundaryless analytic or tropical space; the class is a
  sheaf and restricts to open subsets\\
what is quantified over
 &all models of \(X\), a directed set with no countable cofinal subfamily
  in general (\Cref{lem:ch21-no-cofinal-sequence})
 &neighbourhoods of a point, and decreasing \emph{sequences}\\
positivity of what
 &of a numerical class on a model, tested against vertical curves
 &of a current, tested against test forms\\
values
 &\([-\infty,+\infty)\), bounded above on \(X^{\An}\)
  (\Cref{prop:ch21-sup-on-skeleton})
 &\([-\infty,+\infty)\), bounded above only on relatively compact
  subsets\\
stability under \(\max\)
 &\Cref{thm:ch21-stability}, via \Cref{prop:ch20-model-max}
 &\Cref{prop:ch15-max-stability}, via smooth regularized maxima\\
stability under decreasing limits
 &along arbitrary \emph{nets}
  (\Cref{thm:ch21-decreasing-stability})
 &along \emph{sequences}, by construction; the net statement is not
  available\\
\bottomrule
\end{tabular}
\end{center}

The last two rows are the crux, and they explain why no formal argument
identifies the classes.  A BFJ psh function is, by
\Cref{thm:bfj-decreasing-regularization}, the limit of a decreasing
\emph{net} of model functions, and \Cref{lem:ch21-no-cofinal-sequence} shows
that this net cannot in general be replaced by a sequence.  The CLD
hierarchy, on the other hand, is generated by decreasing sequences, and
\Cref{thm:ch15-approachable-stability} explicitly records that stability
under further decreasing limits is exactly what is not known.  So the
statement ``both are psh, hence they agree'' is not an argument in either
direction; it is a pun.

What the book does prove is
\Cref{thm:ch19-cld-bfj-comparison}: on the model locus the two notions
coincide, together with Zhang semipositivity, positivity of the Chern
current, positivity of the Chern delta-form, functorial pshness of the local
weights, and the curve test; and on the continuous locus the two associated
Monge--Amp\`ere measures coincide.  Outside those two loci the book proves
nothing, and \Cref{thm:ch19-cld-bfj-comparison}\,(C) lists the four
statements that are \emph{not} being asserted.  \Cref{sec:ch21-cld-contrast}
returns to this with the two provable structural separations and with a
precise account of why no separating example is available in the present
scope.

\subsection{What is proved here, what is imported, what is deferred}
\label{sec:ch21-imported}

Proved here in full: the definition and its cofinality
(\Cref{sec:ch21-definition}); the model-consistency theorem
(\Cref{thm:ch21-model-consistency}), granted the multiplier-ideal package
named in \Cref{sec:ch21-multiplier-input}; the equivalent skeletal
formulations (\Cref{thm:ch21-equivalent}); the reconstruction theorem
(\Cref{thm:bfj-skeletal-characterization}); the entire stability calculus
(\Cref{sec:ch21-stability}); the polar-set theory
(\Cref{sec:ch21-singularities}); the curve calculus and all the worked
examples.

Imported with exact citations: the asymptotic multiplier-ideal package of
\cite[Appendix~B]{BFJ16}, restated as
\eqref{eq:ch21-multiplier-subadditivity}--\eqref{eq:ch21-uniform-generation}
and used only in \Cref{thm:ch21-model-consistency}; the skeletal
equicontinuity estimate \Cref{thm:bfj-skeletal-equicontinuity}, proved in
\Cref{chap:bfj-compactness} and used only in
\Cref{thm:ch21-decreasing-stability}; and the decreasing regularization
theorem \Cref{thm:bfj-decreasing-regularization}, proved in
\Cref{chap:bfj-compactness} from the envelope regularity of
\Cref{chap:bfj-envelopes}, used only in the continuous comparison
\Cref{thm:ch21-continuous-zhang}.  All three are stated before use, and none
is used to prove any statement on which they themselves depend; the
dependency chain is displayed in \Cref{sec:ch21-ledger}.

Deferred: the Monge--Amp\`ere operator (\Cref{chap:bfj-ma}), capacity and
quasicontinuity (\Cref{chap:bfj-capacity}), envelopes
(\Cref{chap:bfj-envelopes}), and pullback beyond the model and continuous
loci (\Cref{chap:bfj-functoriality}).  In particular no measure is attached
to an unbounded psh function anywhere in this chapter, and the reader should
resist reading \(\theta+\ddc\varphi\) as a measure before
\Cref{chap:bfj-ma} has constructed one.

\subsection{Plan}
\label{sec:ch21-plan}

\Cref{sec:ch21-definition} gives the definition, shows that neither clause
can be dropped, and proves cofinality independence.
\Cref{sec:ch21-consistency} proves that the new notion agrees with model
semipositivity on model functions; this is the one deep theorem of the
chapter.  \Cref{sec:ch21-equivalent} assembles the equivalent skeletal
formulations, and \Cref{sec:ch21-restriction} proves the reconstruction
theorem and studies the inverse system of restrictions to dual complexes.
\Cref{sec:ch21-stability} is the stability calculus.
\Cref{sec:ch21-singularities} treats singularity types and polar sets.
\Cref{sec:ch21-continuous} compares with continuous semipositive metrics and
delimits the comparison with the CLD class.  \Cref{sec:ch21-worked} carries seven examples to explicit numbers, and
\Cref{sec:ch21-dictionary} is the chapter dictionary.

\section{The definition}
\label{sec:ch21-definition}

\subsection{Skeleta and retractions: notation}
\label{sec:ch21-notation}

If \(\mathcal X\) is an SNC model, write

\begin{equation}
\label{eq:ch21-skeletal-maps}
 \iota_{\mathcal X}:\Delta_{\mathcal X}\hookrightarrow X^{\An},
 \qquad
 p_{\mathcal X}:X^{\An}\longrightarrow
       \iota_{\mathcal X}(\Delta_{\mathcal X})
\end{equation}

for its embedding and retraction.  We routinely identify the dual complex
with its image.

The notation of \Cref{chap:skeleta} is
\(\Sk(\mathcal X)=\iota_{\mathcal X}(\Delta_{\mathcal X})\) for the image
and \(r_{\mathcal X}\) for the retraction; the two systems agree under the
identification just made, so that \(p_{\mathcal X}=r_{\mathcal X}\) and
\(\Delta_{\mathcal X}=\Sk(\mathcal X)\) whenever both appear.  We use
\(\Delta_{\mathcal X}\) here because the polyhedral structure, and not only
the underlying subset of \(X^{\An}\), is in constant use: every statement
below about convexity refers to the integral affine structure of
\Cref{def:integral-affine}, transported by \(\iota_{\mathcal X}\).

Four facts from \Cref{chap:skeleta} are used without further comment.

\begin{enumerate}
\item \emph{Nesting.}  If \(\mathcal X'\geq\mathcal X\) are SNC models, then
\(\Delta_{\mathcal X}\subset\Delta_{\mathcal X'}\) and the cells of
\(\Delta_{\mathcal X}\) are unions of cells of \(\Delta_{\mathcal X'}\)
inducing the same integral affine structure
(\Cref{prop:skeleton-nesting}, \Cref{prop:affine-structure-refinement}).
The retractions are then ordered by domination:

\begin{equation}
\label{eq:ch21-retraction-composition}
 p_{\mathcal X}p_{\mathcal X'}=p_{\mathcal X},
 \qquad
 p_{\mathcal X'}p_{\mathcal X}=p_{\mathcal X}.
\end{equation}

The asymmetry of these two identities reflects the inclusion of the old
skeleton into the new one: the first says that retracting twice is
retracting once, the second that the coarser skeleton is already fixed by
the finer retraction.
\item \emph{Cofinality.}  SNC models are cofinal among models and form a
directed set (\Cref{prop:snc-cofinal}); this is where resolution of
singularities, hence \(\operatorname{char}\widetilde K=0\), enters.
\item \emph{Inverse limit.}  \(X^{\An}=\varprojlim_{\mathcal X}
\Delta_{\mathcal X}\) as a topological space, so
\(p_{\mathcal X}(x)\to x\) along the directed set of models, and
\(\bigcup_{\mathcal X}\Delta_{\mathcal X}\) is dense in \(X^{\An}\)
(\Cref{thm:skeleton-inverse-limit}).
\item \emph{Divisorial normalization.}  The vertex of
\(\Delta_{\mathcal X}\) attached to a component \(E\) of multiplicity
\(b_E\) is \(x_E=b_E^{-1}\ord_E\), not \(\ord_E\)
(\Cref{rem:divisorial-normalization}).  Every numerical statement below
depends on this.
\end{enumerate}

\subsection{The two clauses}
\label{sec:ch21-two-clauses}

\begin{definition}\label{def:ch21-bfj-psh}
A function

\[
             \varphi:X^{\An}\longrightarrow[-\infty,+\infty)
\]

is \emph{\(\theta\)-psh} if it is upper semicontinuous, not identically
\(-\infty\), and, for every SNC model \(\mathcal X\) on which \(\theta\) is
determined, the following hold:

\begin{enumerate}
\item
\begin{equation}
\label{eq:ch21-retraction-inequality}
                    \varphi\leq\varphi\circ p_{\mathcal X}
                    \quad\text{on }X^{\An};
\end{equation}

\item the restriction \(\varphi|_{\Delta_{\mathcal X}}\) is a uniform limit
of restrictions of \(\theta\)-psh model functions.
\end{enumerate}

The resulting class is denoted \(\PSH(X,\theta)\).
\end{definition}

Four remarks on the wording are in order before anything is proved.

\emph{The quantifier in clause~(2).}  The model functions occurring there
are not required to be determined on \(\mathcal X\); they range over
\emph{all} \(\theta\)-psh model functions, hence over all models.  Only
their restrictions to the fixed complex \(\Delta_{\mathcal X}\) are compared
with \(\varphi\).  This is essential: a \(\theta\)-psh model function
determined on \(\mathcal X\) is affine on each face of
\(\Delta_{\mathcal X}\) by \Cref{lem:model-function-affine}, and the affine
ones are far too rigid to approximate an arbitrary convex function on a
simplex.  What is being asked is that
\(\varphi|_{\Delta_{\mathcal X}}\) lie in the uniform closure

\begin{equation}
\label{eq:ch21-skeletal-closure}
 \mathcal P_{\mathcal X}(\theta)
 :=\overline{
   \bigl\{\psi|_{\Delta_{\mathcal X}}:
   \psi\ \text{a \(\theta\)-psh model function}\bigr\}
   }^{\,\lVert\cdot\rVert_\infty}
   \ \subset\ C^0(\Delta_{\mathcal X}).
\end{equation}

\emph{The uniform norm.}  The closure in
\eqref{eq:ch21-skeletal-closure} is taken for the sup-norm on the compact
polyhedron \(\Delta_{\mathcal X}\), not for pointwise convergence.  Pointwise
convergence would be a strictly weaker requirement and would not force
continuity of the restriction; uniform convergence does, which is what makes
\Cref{thm:bfj-skeletal-characterization}\,(1) available at once.  On the
other hand \Cref{thm:ch21-equivalent} shows that pointwise convergence at
divisorial points may be substituted, provided the retraction inequality is
retained --- the missing compactness being supplied by the equicontinuity of
\Cref{thm:bfj-skeletal-equicontinuity}.

\emph{Upper semicontinuity and the value \(-\infty\).}  Both are forced by
the intended applications: the regularized supremum of a family of
\(\theta\)-psh functions is upper semicontinuous but generally not
continuous, and a decreasing limit of \(\theta\)-psh functions generally
takes the value \(-\infty\) somewhere.  The exclusion of
\(\varphi\equiv-\infty\) is a convention, not a theorem; without it the
class would fail to be a subset of the space in which
\Cref{chap:bfj-energy} works, because no reference point of finiteness would
be available.  The extended-real conventions of
\Cref{sec:extended-arithmetic} are in force: no difference of two
\(-\infty\)'s is ever formed, and every identity below is either
monotone or asserted only where both sides are defined.

\emph{Dependence on \(\theta\).}  Membership in \(\PSH(X,\theta)\) is not a
property of \(\varphi\) alone.  It is a property of the pair
\((\theta,\varphi)\), equivalently of the singular metric
\(\|\cdot\|_0\rme^{-\varphi}\) once a model reference of curvature
\(\theta\) has been chosen; this is the content of
\Cref{cor:ch21-reference-independence}.  The conventions of the book
(\Cref{sec:appJ-model-functions} and the paragraph on positivity words)
forbid inferring anything about a psh class from its notation alone, and
\Cref{sec:ch21-global-vs-local} is the reason for that prohibition.

Clause (2) implies immediately that the restriction to each dual complex is
finite and continuous.  It is convex on each face because model psh
restrictions are convex and uniform limits preserve convexity.  Notice that
this does not force \(\varphi\) to be finite on all of \(X^{\An}\).

\begin{proposition}
\label{prop:ch21-first-consequences}
Let \(\varphi\in\PSH(X,\theta)\) and let \(\mathcal X\) be an SNC model on
which \(\theta\) is determined.  Then:
\begin{enumerate}
\item\label{itm:ch21-fc:finite}
\(\varphi|_{\Delta_{\mathcal X}}\) is real valued and continuous, and it is
convex on every face of \(\Delta_{\mathcal X}\) for the integral affine
structure of \Cref{def:integral-affine};
\item\label{itm:ch21-fc:bounded}
\(\varphi\) is bounded above on \(X^{\An}\), with
\(\sup_{X^{\An}}\varphi=\max_{\Delta_{\mathcal X}}\varphi\);
\item\label{itm:ch21-fc:quasimonomial}
\(\varphi\) is finite at every quasimonomial point, so that
\(\{\varphi=-\infty\}\) is disjoint from
\(\bigcup_{\mathcal Y}\Delta_{\mathcal Y}\);
\item\label{itm:ch21-fc:emptyinterior}
\(\{\varphi=-\infty\}\) has empty interior in \(X^{\An}\).
\end{enumerate}
\end{proposition}

\begin{proof}
\cref{itm:ch21-fc:finite}  Let \(\psi_k\) be \(\theta\)-psh model functions
with \(\psi_k|_{\Delta_{\mathcal X}}\to\varphi|_{\Delta_{\mathcal X}}\)
uniformly.  Each \(\psi_k\) is a real-valued continuous function on
\(X^{\An}\) by \eqref{eq:ch20-model-functions} and the embedding of model
functions into \(C^0(X^{\An})\); a uniform limit of real-valued continuous
functions on a compact space is real valued and continuous.  Convexity on
faces holds for each \(\psi_k\) by
\Cref{lem:ch21-model-skeletal-properties}, and convexity is preserved by
pointwise limits, in particular by uniform ones.

\cref{itm:ch21-fc:bounded}  This is \Cref{prop:ch21-sup-on-skeleton}, proved
below from \eqref{eq:ch21-retraction-inequality} alone; we record it here
because it is used constantly.  Note that it is the retraction inequality,
and not compactness of \(X^{\An}\), that gives the sharper statement that
the supremum is attained on \emph{every} determining dual complex.

\cref{itm:ch21-fc:quasimonomial}  A quasimonomial point lies on
\(\Delta_{\mathcal Y}\) for some SNC model \(\mathcal Y\); enlarging
\(\mathcal Y\) so that it also determines \(\theta\) preserves this by the
nesting property.  Now apply \cref{itm:ch21-fc:finite} to \(\mathcal Y\).

\cref{itm:ch21-fc:emptyinterior}  Let \(U\subset X^{\An}\) be nonempty and
open.  By \Cref{thm:skeleton-inverse-limit} the union of the skeleta is
dense, so \(U\) meets some \(\Delta_{\mathcal Y}\), and
\cref{itm:ch21-fc:quasimonomial} gives a point of \(U\) at which \(\varphi\)
is finite.
\end{proof}

\subsection{Neither clause may be dropped}
\label{sec:ch21-clauses-independent}

Clause (1) cannot be omitted.  A compatible-looking collection of convex
functions on some selected polyhedra need not prescribe the value at an
analytic point outside their union.  Inequality
\eqref{eq:ch21-retraction-inequality}, together with upper
semicontinuity and convergence of the retractions, supplies the missing
global information.  The following example makes this completely explicit
and shows that the failure is not a matter of pathology in the value at one
point: the function exhibited is upper semicontinuous, satisfies clause~(2)
on every model, and is not \(\theta\)-psh.

\begin{example}[Clause~(2) alone is not enough]
\label{ex:ch21-clause-one-needed}
\index{plurisubharmonic function!BFJ}
Assume \eqref{eq:ch21-hypotheses} and let \(\theta\) be semipositive, so
that the constant \(0\) is a \(\theta\)-psh model function.  Choose a point
\(x_0\in X^{\An}\) that is not quasimonomial; for instance a
\(K\)-rational point of \(X\), viewed in \(X^{\An}\).  It lies on no
skeleton because it has nonzero schematic support, whereas every point of a
model skeleton is a valuation of the function field \(K(X)\), hence has
zero support
(\Cref{prop:berk-pair-description,lem:skeleton-monomial-invariants}).
Define
\[
 \varphi:=\one_{\{x_0\}},
 \qquad\text{that is,}\quad
 \varphi(x_0)=1,\quad \varphi(x)=0\ \text{ for }x\neq x_0 .
\]

Then \(\varphi\) is upper semicontinuous.  Indeed \(\{\varphi<c\}\) is empty
for \(c\leq0\), equals \(X^{\An}\setminus\{x_0\}\) for \(0<c\leq1\), and
equals \(X^{\An}\) for \(c>1\); the middle set is open because \(X^{\An}\) is
Hausdorff (\Cref{thm:berk-proper-compact}), so points are closed.  It is not
identically \(-\infty\).  Its restriction to every dual complex is
identically \(0\), because \(x_0\) lies on no dual complex; hence
clause~(2) holds for every SNC model, with the constant approximating
sequence \(\psi_k\equiv0\).

Clause~(1) fails at \(x_0\): for every SNC model,
\(\varphi(p_{\mathcal X}(x_0))=0<1=\varphi(x_0)\).  And \(\varphi\notin
\PSH(X,\theta)\): if it were, \Cref{thm:bfj-skeletal-characterization}
would give \(\varphi(x_0)=\inf_{\mathcal X}\varphi(p_{\mathcal X}(x_0))=0\).

Two features of this example are worth recording.  First, the defect is
invisible on every dual complex, so no amount of skeletal information
detects it; the retraction inequality is a genuinely extra axiom, not a
normalization.  Second, the example uses only that some point of
\(X^{\An}\) is not quasimonomial, which is true for every positive
dimensional \(X\): already on a curve the type-\(1\) and type-\(4\) points
lie on no skeleton (\Cref{thm:berk-a1-classification},
\Cref{rem:quasimonomial-vs-divisorial}\,(3)).
\end{example}

Clause (2) cannot be omitted either, and for a different reason: without it
nothing at all remembers positivity.

\begin{example}[Clause~(1) alone is not enough]
\label{ex:ch21-clause-two-needed}
Keep the notation of \Cref{ex:ch21-clause-one-needed} and let \(\mathcal X\)
be an SNC model determining \(\theta\).  Let \(h:\Delta_{\mathcal X}\to
\mathbf R\) be any continuous function that is \emph{concave} and
nonconstant on some face --- for instance, in the curve case of
\Cref{ex:ch21-conic-numbers}, \(h(t)=-|t-\tfrac12|\) on
\(\Delta_{\mathcal X}=[0,1]\).  Put \(\varphi:=h\circ p_{\mathcal X}\).

Then \(\varphi\) is continuous, hence upper semicontinuous, and finite; and
clause~(1) holds, with equality, for every SNC model
\(\mathcal Y\geq\mathcal X\), a family which is cofinal in the models
determining \(\theta\).  Indeed
\(p_{\mathcal X}\circ p_{\mathcal Y}=p_{\mathcal X}\) by
\eqref{eq:ch21-retraction-composition}, so \(\varphi\circ p_{\mathcal Y}=
h\circ p_{\mathcal X}\circ p_{\mathcal Y}=h\circ p_{\mathcal X}=\varphi\),
and equality is a case of the inequality.  But clause~(2) fails at
\(\mathcal X\): every element of \(\mathcal P_{\mathcal X}(\theta)\) is
convex on each face by
\Cref{lem:ch21-model-skeletal-properties} and the stability of convexity
under uniform limits, whereas
\(\varphi|_{\Delta_{\mathcal X}}=h\) is strictly concave along a segment
inside a face.  Hence \(\varphi\notin\PSH(X,\theta)\).

The two examples together show that the clauses are logically independent:
neither implies the other, and each excludes functions that the other
admits.  It is worth noting where cofinality enters.  In
\Cref{ex:ch21-clause-one-needed} clause~(2) holds for \emph{every} model and
clause~(1) fails; here clause~(1) holds on a cofinal family and clause~(2)
fails.  One cannot do better in the second case without extra work, because
the proof of \Cref{prop:ch21-cofinal-independence} propagates clause~(1)
from a cofinal family to all models by using clause~(2), which is exactly
what is being denied.
\end{example}

\subsection{Cofinality}
\label{sec:ch21-cofinality}

Verifying \Cref{def:ch21-bfj-psh} for \emph{every} SNC model determining
\(\theta\) would be unusable in practice.  It is not necessary.

\begin{proposition}\label{prop:ch21-cofinal-independence}
It is enough to verify \Cref{def:ch21-bfj-psh} on any cofinal directed family
of SNC models on which \(\theta\) is determined.
\end{proposition}

\begin{proof}
Let \(\mathfrak M\) be the given cofinal directed family and suppose the two
clauses hold for every member of \(\mathfrak M\).  Let \(\mathcal Y\) be an
arbitrary SNC model on which \(\theta\) is determined and choose
\(\mathcal X\in\mathfrak M\) dominating \(\mathcal Y\); such an
\(\mathcal X\) exists by cofinality.  By the nesting property recalled in
\Cref{sec:ch21-notation}, \(\Delta_{\mathcal Y}\subset\Delta_{\mathcal X}\).

\emph{Clause~(2) for \(\mathcal Y\).}  Let \(\psi_k\) be \(\theta\)-psh
model functions with
\(\psi_k|_{\Delta_{\mathcal X}}\to\varphi|_{\Delta_{\mathcal X}}\)
uniformly.  Restriction to a subset does not increase the sup-norm, so
\(\psi_k|_{\Delta_{\mathcal Y}}\to\varphi|_{\Delta_{\mathcal Y}}\)
uniformly.  This is clause~(2) for \(\mathcal Y\).

\emph{Clause~(1) for \(\mathcal Y\).}  By
\Cref{lem:ch21-model-skeletal-properties} applied to \(\mathcal Y\), each
approximant satisfies \(\psi_k\leq\psi_k\circ p_{\mathcal Y}\).  Fix
\(x\in X^{\An}\).  Both \(p_{\mathcal X}(x)\) and
\(p_{\mathcal Y}p_{\mathcal X}(x)=p_{\mathcal Y}(x)\) lie in
\(\Delta_{\mathcal X}\), the second because
\(\Delta_{\mathcal Y}\subset\Delta_{\mathcal X}\) and because
\(p_{\mathcal Y}\circ p_{\mathcal X}=p_{\mathcal Y}\) by
\eqref{eq:ch21-retraction-composition}.  Passing to the uniform limit of the
inequality \(\psi_k\leq\psi_k\circ p_{\mathcal Y}\) at the point
\(p_{\mathcal X}(x)\), where the convergence is uniform, gives

\begin{equation}
\label{eq:ch21-cofinal-retractions}
 \varphi(p_{\mathcal X}(x))
 \leq
 \varphi(p_{\mathcal Y}p_{\mathcal X}(x))
 =\varphi(p_{\mathcal Y}(x)).
\end{equation}

Together with \(\varphi(x)\leq\varphi(p_{\mathcal X}(x))\), which is
clause~(1) for \(\mathcal X\in\mathfrak M\), this proves
clause~(1) for \(\mathcal Y\).
\end{proof}

\begin{remark}[Which cofinal families are used]
\label{rem:ch21-cofinal-families}
Three cofinal families recur.  \emph{(a)} The SNC models dominating one
fixed model \(\mathcal X_0\) determining \(\theta\); this is the default and
is directed because two models are dominated by a third
(\Cref{prop:snc-cofinal}\,(2)).  \emph{(b)} The vertically
\(\mathbf Q\)-factorial SNC models, used in
\Cref{chap:bfj-compactness}.  \emph{(c)} The toroidal blowups of a fixed
\(\mathcal X_0\), which do not enlarge \(\Delta_{\mathcal X_0}\) as a subset
of \(X^{\An}\) but subdivide it (\Cref{thm:skeleton-subdivision}); this
family is \emph{not} cofinal, and using it alone would test only the
behaviour of \(\varphi\) on one skeleton.  The distinction between (a) and
(c) is exactly the distinction between enlarging a skeleton and refining
its cell structure, and confusing them is the most common error in
skeletal arguments.
\end{remark}

\subsection{Why nets and not sequences}
\label{sec:ch21-nets}

\index{net!no cofinal sequence of models}
Every limit taken in this chapter and the next is indexed by the directed
set of models, and the book's standing convention
(\Cref{rem:net-vs-sequence-checklist}) forbids replacing such a limit by a
sequential one without a proof.  The following lemma shows that the
prohibition is not pedantry: already for the projective line over a very
ordinary discretely valued field of residue characteristic zero, no
countable family of models is cofinal.  Its proof is the model-theoretic
counterpart of \Cref{ex:no-cofinal-sequence}.

\begin{lemma}
\label{lem:ch21-no-cofinal-sequence}
Let \(K\) be complete and discretely valued with uncountable residue field
\(\widetilde K\); for instance \(K=\mathbf C((t))\), which satisfies all of
\eqref{eq:ch21-hypotheses}.  Let \(X=\mathbf P^1_K\).  Then the directed set
of SNC models of \(X\) over \(R\), ordered by domination, has no countable
cofinal subset.  Consequently no decreasing net indexed by that set may be
replaced by a decreasing sequence merely because the index set is directed.
\end{lemma}

\begin{proof}
For \(a\in R\) let \(\zeta_a:=\zeta_{a,|\pi|}\in\mathbf P^{1,\An}_K\) be the
Gauss point of the closed disc of centre \(a\) and radius \(|\pi|\), in the
notation of \Cref{sec:berkovich-p1}.  Two such points coincide exactly when
the discs coincide, that is, when \(|a-b|\leq|\pi|\), that is, when \(a\) and
\(b\) have the same image in \(R/\pi R\).  Since
\(R/\pi R\cong\widetilde K\) is uncountable, the family
\(\{\zeta_a\}\) contains uncountably many distinct points.  Each is
divisorial: it is the divisorial point of the exceptional component of the
blowup of \(\mathbf P^1_R\) at the closed point \(\bar a\) of the special
fibre, a component of multiplicity one, and its radius \(|\pi|\) lies in
\(|K^\times|\).  Hence each \(\zeta_a\) lies on the dual complex of some SNC
model, by \Cref{rem:quasimonomial-vs-divisorial}\,(1).

Each \(\zeta_a\) lies at distance \(1\) from the Gauss point
\(\eta=\zeta_{0,1}\) in the path metric of
\Cref{def:ch08-path-metric}, and for \(a\not\equiv b\) the paths
\([\eta,\zeta_a]\) and \([\eta,\zeta_b]\) meet only in \(\eta\).  A dual
complex of an SNC model of a curve is a finite graph, hence a finite union
of segments, and therefore contains only finitely many points \(\zeta_a\):
each such point would require its own branch at \(\eta\), and a finite graph
has finitely many branches at a point.

Now suppose \((\mathcal X_j)_{j\geq1}\) were a countable cofinal family.
Then \(\bigcup_j\Delta_{\mathcal X_j}\) contains only countably many of the
\(\zeta_a\), so there is \(a\) with
\(\zeta_a\notin\bigcup_j\Delta_{\mathcal X_j}\).  Choose an SNC model
\(\mathcal Y\) with \(\zeta_a\in\Delta_{\mathcal Y}\).  By cofinality some
\(\mathcal X_j\) dominates \(\mathcal Y\), whence
\(\Delta_{\mathcal Y}\subset\Delta_{\mathcal X_j}\) by the nesting property,
and \(\zeta_a\in\Delta_{\mathcal X_j}\); a contradiction.

The final sentence is a formality: if a decreasing net
\((u_i)_{i\in I}\) were replaced by a sequence \((u_{i_k})\) with the same
infimum, the indices \(i_k\) would have to be cofinal in \(I\) whenever the
net separates points, which the preceding paragraphs forbid.
\end{proof}

\begin{remark}[Scope of the lemma]
\label{rem:ch21-nets-scope}
The lemma is a statement about the \emph{index set}, not about any
particular psh function, and it is not a claim that sequential arguments are
never available.  Three legitimate sequential reductions are used in this
part, all listed in \Cref{rem:net-vs-sequence-checklist}: a single dual
complex is a compact metrizable space, so sequences suffice on it
(\Cref{prop:ch23-fixed-skeleton-selection}); the Choquet selection lemma
\Cref{lem:choquet-selection} replaces a family by a countable subfamily with
the same regularized supremum whenever the ambient space is second
countable; and a concrete tower of blowups, such as the one in
\Cref{ex:ch21-p1-refinement}, is indexed by \(\mathbf N\) by construction.
What is never legitimate is to pass from the full model system to a
sequence without naming one of these three mechanisms.
\end{remark}

\section{Consistency with model semipositivity}
\label{sec:ch21-consistency}

The word ``psh'' has now appeared twice: first as relative nefness for a
model function, and then in \Cref{def:ch21-bfj-psh}.  Their agreement is not
formal.

One direction is elementary and is the content of
\Cref{lem:ch21-model-skeletal-properties}: a relatively nef vertical divisor
produces a function that is convex on faces and dominated by its own
retraction, so it satisfies both clauses.  The other direction is a
closedness statement, and it is the one deep theorem of the chapter.  To see
why it cannot be formal, notice what it asserts.  Clause~(2) says that the
\emph{restriction} of \(\varphi\) to each dual complex is approximable;
relative nefness is a statement about degrees on vertical curves inside a
model, and a vertical curve of an \(n\)-dimensional special fibre is not
visible on the dual complex at all when \(n\geq2\).  There is therefore a
genuine gap between skeletal data and numerical positivity, and the theorem
below closes it by an argument that runs through base ideals, asymptotic
multiplier ideals, and uniform global generation.

\subsection{Skeletal properties of model psh functions}
\label{sec:ch21-model-skeletal}

\begin{lemma}\label{lem:ch21-model-skeletal-properties}
Let \(\mathcal X\) be an SNC determination of \(\theta\).  Every
\(\theta\)-psh model function \(\varphi\) is convex on each face of
\(\Delta_{\mathcal X}\) and satisfies

\[
                 \varphi\leq\varphi\circ p_{\mathcal X}.
\]

If \(\varphi\) is determined on \(\mathcal X\), equality holds in the last
display.
\end{lemma}

\begin{proof}
\emph{Step 1: reduction to an ideal.}  Assume first that
\(\varphi=\varphi_D\) is rational, and choose a normalized
vertical blowup
\(\rho:\mathcal Y\to\mathcal X\) on which \(D\) is Cartier.  Because
\(\theta\) comes from \(\mathcal X\), its pullback has degree zero on every
curve contracted by \(\rho\), by the first case of
\eqref{eq:ch20-nef-pullback-test}.  The nefness of
\(\rho^*\theta_{\mathcal X}+[D]\) therefore says that \(D\) is
\(\rho\)-nef: for a curve \(C\) contracted by \(\rho\),
\((D\cdot C)=(\rho^*\theta_{\mathcal X}+[D])\cdot C\geq0\).  A relatively
nef vertical divisor is a limit of relatively
ample vertical \(\mathbf Q\)-divisors \(A_k\), by Kleiman's criterion
applied in the finite-dimensional space of vertical numerical classes on
\(\mathcal Y\) relative to \(\mathcal X\).  After
replacing \(A_k\) by a
multiple, \(\mathcal O_{\mathcal Y}(m_kA_k)\) is generated relative to
\(\mathcal X\).  With

\[
       \mathfrak a_k:=\rho_*
          \mathcal O_{\mathcal Y}(m_kA_k),
\]

the normalized blowup identity gives
\(m_k^{-1}\log|\mathfrak a_k|=\varphi_{A_k}\), and these functions converge
uniformly to \(\varphi_D\).  This is the relative Kleiman approximation used
in \cite[Lemma~5.3]{BFJ16}.

\emph{Step 2: convexity of \(\log|\mathfrak a_k|\) on a face.}  Locally on
\(\mathcal X\), write \(\mathfrak a_k=(f_1,\ldots,f_N)\).  On a face
\(\sigma_J\) of the dual complex, the point \(\iota_{\mathcal X}(\alpha)\)
is the monomial valuation \(v_\alpha\), and
\Cref{lem:monomial-newton} evaluates \(v_\alpha(f_j)\) as the minimum of the
finitely many linear functions \(\alpha\mapsto\langle\alpha,\beta\rangle\)
over the exponents \(\beta\) occurring in the expansion of \(f_j\).  Hence
\(\log|f_j|=-v_\bullet(f_j)\) is a maximum of finitely many affine
functions on \(\sigma_J\), and

\[
       \log|\mathfrak a_k|
       =\max_{1\leq j\leq N}\log|f_j|
\]

is again such a maximum, hence convex.  This is
\Cref{lem:vertical-ideal-convex}\,(2), whose proof is exactly this
computation; we have repeated it because the convexity, and not merely the
piecewise affineness, is what passes to the limit.

\emph{Step 3: the retraction inequality for \(\log|\mathfrak a_k|\).}  The
Gauss property of the monomial valuation defining the retraction gives
\(|f_j(x)|\leq|f_j(p_{\mathcal X}(x))|\) for every \(j\); this is
\Cref{lem:retraction-minorant}, applied at the center of \(x\).  Taking the
maximum over \(j\) and using
\Cref{cor:retraction-ideal-monotone} yields

\[
 \log|\mathfrak a_k|(x)
 \leq\log|\mathfrak a_k|(p_{\mathcal X}(x)).
\]

\emph{Step 4: passage to the limit and the real case.}  Uniform passage to
the limit preserves both a facewise convexity inequality
\(u(\tfrac12(\alpha+\alpha'))\leq\tfrac12(u(\alpha)+u(\alpha'))\) and the
pointwise inequality of Step~3, so both assertions hold for rational
\(\varphi\).  A real model function is obtained as a limit, on one common
model, of rational vertical divisors after an arbitrarily small ample
perturbation of the curvature class, exactly as in the last paragraph of the
proof of \Cref{prop:ch20-model-max}; closedness of the nef cone and uniform
convergence of divisor functions give the same conclusions.

\emph{Step 5: equality on a determining model.}  Finally, if
\(D\) already lies on \(\mathcal X\), then \(\varphi_D\) is affine on each
face of \(\Delta_{\mathcal X}\) and satisfies
\(\varphi_D\circ p_{\mathcal X}=\varphi_D\) by
\Cref{lem:model-function-affine}; that is, a model function determined on
\(\mathcal X\) is constant on the fibres of \(p_{\mathcal X}\), so
\(\varphi_D=\varphi_D\circ p_{\mathcal X}\).
\end{proof}

\begin{remark}[The three ingredients, separated]
\label{rem:ch21-lemma-ingredients}
It is worth isolating what each step of the last proof contributes, because
each is used again independently.  Step~1 is positivity: it converts
relative nefness into relative global generation, at the cost of an
approximation.  Steps~2 and~3 are the skeletal calculus of
\Cref{chap:skeleta} and use no positivity whatever: \emph{every} vertical
ideal has a convex logarithm on faces and satisfies the retraction
inequality (\Cref{lem:vertical-ideal-convex}).  Step~5 is bookkeeping.  The
only place where positivity is needed is therefore the passage from a
divisor to an ideal, and this is precisely the mechanism that the maximum
construction of \Cref{prop:ch20-model-max} also uses.  A reader who
remembers one thing from this proof should remember that convexity on a face
comes from \emph{ideals}, not from nefness.
\end{remark}

\subsection{The imported multiplier-ideal package}
\label{sec:ch21-multiplier-input}

The converse below uses the following proved package, in exactly the form of
\cite[Appendix~B]{BFJ16}, and reviewed with its proof from log resolutions,
relative Kawamata--Viehweg vanishing and Castelnuovo--Mumford regularity in
\Cref{sec:appC-asymptotic-multiplier}.  Let \(\mathfrak a_\bullet\) be a
graded sequence of vertical ideals on an SNC model \(\mathcal X\), so that
\(\mathfrak a_m\mathfrak a_\ell\subset\mathfrak a_{m+\ell}\), and let
\(\mathcal J(\mathfrak a_\bullet^{\,c})\) denote its asymptotic multiplier
ideal of exponent \(c\).  Then

\begin{equation}
\label{eq:ch21-multiplier-subadditivity}
 \mathfrak a_m\subset
 \mathcal J(\mathfrak a_\bullet^m),
 \qquad
 \mathcal J(\mathfrak a_\bullet^{m\ell})
 \subset\mathcal J(\mathfrak a_\bullet^m)^\ell
 \qquad(m,\ell\geq1),
\end{equation}

and there is an ample line bundle \(\mathcal A\) on \(\mathcal X\),
\emph{independent of \(m\)}, such that

\begin{equation}
\label{eq:ch21-uniform-generation}
 m\mathcal L+\mathcal A\ \text{is globally generated whenever}\
 \mathcal J(\mathfrak a_\bullet^m)=\mathcal O_{\mathcal X},
\end{equation}

\(\mathfrak a_m\) being the base ideal of \(m\mathcal L\).  The precise
statement behind \eqref{eq:ch21-uniform-generation} is that the adjoint
twist \(m\mathcal L+\mathcal A\otimes\mathcal J(\mathfrak a_\bullet^m)\) is
globally generated for one fixed ample \(\mathcal A\) and every \(m\); when
the multiplier ideal is trivial this is the displayed assertion.

Three comments delimit the import.  First,
\eqref{eq:ch21-multiplier-subadditivity} and
\eqref{eq:ch21-uniform-generation} are the \emph{only} facts imported into
the proof of \Cref{thm:ch21-model-consistency}; every deduction from them is
written out below.  Second, both use
\(\operatorname{char}\widetilde K=0\), through the existence of log
resolutions and through Kawamata--Viehweg vanishing; this is the unique
appearance of that hypothesis in the present chapter, and
\Cref{chap:bfj-compactness} discusses what is known without it.  Third,
multiplier ideals are integrally closed, a fact used once below through the
valuative criterion: an integrally closed ideal is determined by the values
\(\log|\mathfrak a|(x)\) at divisorial points \(x\).

\begin{theorem}\label{thm:ch21-model-consistency}
For a model function \(\varphi\), the following are equivalent:

\begin{enumerate}
\item \(\theta+\ddc\varphi\) is a semipositive model form;
\item \(\varphi\in\PSH(X,\theta)\) in the sense of
\Cref{def:ch21-bfj-psh}.
\end{enumerate}
\end{theorem}

\begin{proof}
\emph{Step 0: the easy implication.}
Assume first that (1) holds.  The retraction inequality and facewise
convexity follow from \Cref{lem:ch21-model-skeletal-properties}.  A model
function is real valued and continuous, hence upper semicontinuous and not
identically \(-\infty\).  If
\(\mathcal X\) is a determining SNC model, take the constant sequence with
value \(\varphi|_{\Delta_{\mathcal X}}\) in clause (2).  Thus
\(\varphi\) satisfies both clauses of \Cref{def:ch21-bfj-psh}.

The rest of the proof establishes the converse, and it does so through a
closedness statement of independent interest, recorded as
\eqref{eq:ch21-model-closedness} below: the \(\theta\)-psh model functions
form a set that is closed under pointwise convergence \emph{at divisorial
points only}.  The plan is:
\begin{enumerate}
\item[(a)] prove a nefness criterion for a single line bundle
\(\mathcal L\) whose zero weight is a divisorial pointwise limit of
semipositive model weights (Steps~1--4);
\item[(b)] deduce \eqref{eq:ch21-model-closedness} by an ample perturbation
(Step~5);
\item[(c)] manufacture, from clause~(2), a net of model psh functions
converging to \(\varphi\) at every divisorial point, and apply
\eqref{eq:ch21-model-closedness} (Step~6).
\end{enumerate}

\emph{Step 1: the base-ideal sequence and its Fekete limit.}
Consider
an SNC model \(\mathcal X\) and a line bundle \(\mathcal L\) whose generic
restriction is ample.  Suppose that the weight zero of its model metric is a
pointwise limit, at every divisorial point, of weights of semipositive model
metrics on the same generic line bundle.  Let \(\mathfrak a_m\) be the base
ideal of \(m\mathcal L\), that is, the image of the evaluation map
\(H^0(\mathcal X,m\mathcal L)\otimes(m\mathcal L)^{-1}\to
\mathcal O_{\mathcal X}\).  It is vertical for \(m\gg1\) because the generic
restriction of \(\mathcal L\) is ample, so that \(mL\) is globally generated
on \(X\) for \(m\gg1\) and the base locus is contained in
\(\mathcal X_s\).  Moreover
\(\mathfrak a_m\mathfrak a_\ell\subset\mathfrak a_{m+\ell}\).  Hence

\begin{equation}
\label{eq:ch21-base-ideal-fekete}
 h_m:=m^{-1}\log|\mathfrak a_m|
 \quad\text{satisfies}\quad
 \lim_m h_m=\sup_m h_m\leq0
\end{equation}

pointwise, with the limit understood along sufficiently divisible indices.
Indeed, \(\log|\mathfrak a_m\mathfrak a_\ell|=
\log|\mathfrak a_m|+\log|\mathfrak a_\ell|\) and the inclusion
\(\mathfrak a_m\mathfrak a_\ell\subset\mathfrak a_{m+\ell}\) reverses to
\(\log|\mathfrak a_{m+\ell}|\geq
\log|\mathfrak a_m|+\log|\mathfrak a_\ell|\); thus
\(m\mapsto\log|\mathfrak a_m|(x)\) is superadditive for every fixed \(x\)
and Fekete's lemma applies pointwise, the supremum being at most \(0\)
because \(\mathfrak a_m\subset\mathcal O_{\mathcal X}\).

\emph{Step 2: the limit vanishes at divisorial points.}
We claim that the limit in \eqref{eq:ch21-base-ideal-fekete} is zero at
every divisorial point \(x\).  Fix a
rational \(\varepsilon>0\).  Pointwise convergence at \(x\) and at the
finitely many component valuations of \(\mathcal X_s\) supplies a
semipositive model weight \(\varphi_D\) on a higher model such that

\begin{equation}
\label{eq:ch21-closedness-finite-values}
       \varphi_D(x)\geq-\varepsilon,
       \qquad
       \varphi_D(x_{E_i})\leq\varepsilon
       \quad\text{for all }E_i\subset\mathcal X_s.
\end{equation}

By \Cref{lem:ch21-model-skeletal-properties}, convexity places the maximum
on each face at a vertex and the retraction inequality then gives
\(\varphi_D\leq\varepsilon\) on \(X^{\An}\).  With our convention
\(\varphi_{\mathcal X_s}=1\), the divisor

\[
                         D':=D-\varepsilon\mathcal X_s
\]

is anti-effective and
\(\varphi_{D'}(x)\geq-2\varepsilon\).  Notice the minus sign: it is forced by
the convention \eqref{eq:ch20-sign}.  Applying
\Cref{prop:ch20-positive-representative} to the semipositive class
\(c_1(\mathcal L)+[D']\), after raising the model if necessary, gives an
anti-effective vertical \(\mathbf Q\)-divisor \(D''\) such that

\begin{equation}
\label{eq:ch21-closedness-ample-perturbation}
        -\varepsilon\leq\varphi_{D''}\leq0,
        \qquad
        \mathcal L+D'+D''\
        \text{ is relatively ample}.
\end{equation}

Choose \(m\) so that its multiple is globally generated.  Since
\(D'+D''\leq0\), comparison with the evaluation map defining
\(\mathfrak a_m\) gives, on the higher model,

\[
       \mathcal O\bigl(m(D'+D'')\bigr)
       \subset \mathfrak a_m\mathcal O.
\]

Evaluating at \(x\) yields

\begin{equation}
\label{eq:ch21-closedness-base-bound}
             -3\varepsilon
             \leq\varphi_{D'+D''}(x)
             \leq h_m(x)\leq0.
\end{equation}

Since \(\varepsilon\) is arbitrary,
\eqref{eq:ch21-base-ideal-fekete} has limit zero at every divisorial point.

Two details of Step~2 deserve to be spelled out, since they are exactly
where the sign conventions could be lost.  The passage from
\eqref{eq:ch21-closedness-finite-values} to \(\varphi_D\leq\varepsilon\)
on all of \(X^{\An}\) uses both halves of
\Cref{lem:ch21-model-skeletal-properties}: convexity on each face bounds
\(\varphi_D\) on \(\Delta_{\mathcal X}\) by its maximum over the vertices,
which is at most \(\varepsilon\) by the second inequality in
\eqref{eq:ch21-closedness-finite-values}, and the retraction inequality then
propagates the bound off the skeleton.  And \(D'=D-\varepsilon\mathcal X_s\)
is anti-effective because \(\varphi_{D'}=\varphi_D-\varepsilon\leq0\) by
\eqref{eq:ch20-sign} and \(\varphi_{\mathcal X_s}=1\); the subtraction, and
not an addition, is what makes it so.

\emph{Step 3: from base ideals to multiplier ideals.}
Put \(\mathcal J_m=\mathcal J(\mathfrak a_\bullet^m)\).  The two
multiplier-ideal inclusions
\eqref{eq:ch21-multiplier-subadditivity} give

\begin{equation}
\label{eq:ch21-closedness-multiplier-bound}
       \mathfrak a_{m\ell}\subset\mathcal J_m^\ell,
       \qquad
       \sup_{\ell}\frac1\ell\log|\mathfrak a_{m\ell}|
       \leq\log|\mathcal J_m|\leq0.
\end{equation}

Here the first inclusion is
\(\mathfrak a_{m\ell}\subset\mathcal J(\mathfrak a_\bullet^{m\ell})\subset
\mathcal J(\mathfrak a_\bullet^{m})^{\ell}=\mathcal J_m^\ell\), the two
steps being the two halves of
\eqref{eq:ch21-multiplier-subadditivity}; the second display follows by
taking \(\log|\cdot|\), which is superadditive in the same direction, and
dividing by \(\ell\).

\emph{Step 4: triviality of \(\mathcal J_m\), and nefness of
\(\mathcal L\).}
The left side of \eqref{eq:ch21-closedness-multiplier-bound}
is zero at every divisorial point by
\eqref{eq:ch21-base-ideal-fekete}, since
\(\ell^{-1}\log|\mathfrak a_{m\ell}|=m\,h_{m\ell}\to0\); hence
\(\log|\mathcal J_m|=0\) on all divisorial points.  The valuative criterion
for integral closure, and the fact that multiplier ideals are integrally
closed, imply \(\mathcal J_m=\mathcal O_{\mathcal X}\): an integrally closed
ideal whose associated function vanishes at every divisorial point is the
unit ideal.  Uniform global
generation \eqref{eq:ch21-uniform-generation}
now makes \(m\mathcal L+\mathcal A\) globally generated, hence
nef, for one fixed \(\mathcal A\) and every \(m\).  Dividing its class by
\(m\) and letting \(m\to\infty\) proves that \(\mathcal L\) is nef, the nef
cone being closed.

\emph{Step 5: the closedness statement.}
Now let a model function \(\chi\) be a pointwise limit on all divisorial
points of \(\theta\)-psh model functions.  Replacing \(\theta\) by
\(\theta+\ddc\chi\) and subtracting \(\chi\), we may suppose
\(\chi=0\).  Choose a determination \(\mathcal X\) of \(\theta\), and ample
line bundles \(\mathcal A_j\) whose classes form a basis of
\(N^1(\mathcal X/R)\).  The generic class of \(\theta\) is nef because it is
the generic class of every semipositive approximant.  We may therefore
choose arbitrarily small \(\varepsilon_j>0\) so that

\[
          \theta_{\mathcal X}+\sum_j\varepsilon_jc_1(\mathcal A_j)
\]

is a rational class represented by a \(\mathbf Q\)-line bundle
\(\mathcal L_{\varepsilon}\) with ample generic restriction.  Adding the
small ample class to each approximating semipositive metric shows that the
zero weight of \(\mathcal L_{\varepsilon}\) is a divisorial pointwise limit
of semipositive model metrics.  The preceding argument makes
\(\mathcal L_{\varepsilon}\) nef.  Letting all \(\varepsilon_j\to0\) and
using closedness of the nef cone proves that \(\theta\) is semipositive.
Undoing the reference change shows:

\begin{equation}
\label{eq:ch21-model-closedness}
 \left.
 \begin{array}{c}
 \chi_\alpha\text{ is a \(\theta\)-psh model function},\\
 \chi_\alpha\to\chi\text{ at every divisorial point},\\
 \chi\text{ is a model function}
 \end{array}
 \right\}
 \quad\Longrightarrow\quad
 \chi\text{ is a \(\theta\)-psh model function}.
\end{equation}

\emph{Step 6: from clause \textup{(2)} to a divisorial net.}
Finally assume (2).  Let \(I\) be the set of pairs \((S,\varepsilon)\) with
\(S\) a finite set of divisorial points and \(\varepsilon>0\), ordered by
\((S,\varepsilon)\leq(S',\varepsilon')\) if and only if \(S\subset S'\) and
\(\varepsilon\geq\varepsilon'\).  This is a directed set: it is the fifth
entry of \Cref{ex:directed-sets-of-the-book}.  Given \((S,\varepsilon)\),
choose an SNC model \(\mathcal X\) determining \(\theta\) whose dual complex
contains \(S\); such a model exists because each point of \(S\) is
divisorial, hence a vertex on some model
(\Cref{rem:quasimonomial-vs-divisorial}\,(1)), and finitely many models are
dominated by one.  Clause~(2) at \(\mathcal X\) supplies a \(\theta\)-psh
model function \(\chi_{(S,\varepsilon)}\) with
\(\lvert\chi_{(S,\varepsilon)}-\varphi\rvert\leq\varepsilon\) on
\(\Delta_{\mathcal X}\), in particular on \(S\).

The resulting net \((\chi_i)_{i\in I}\) converges to \(\varphi\) at every
divisorial point \(x\): given \(\delta>0\), every
\(i\geq(\{x\},\delta)\) satisfies \(|\chi_i(x)-\varphi(x)|\leq\delta\).
Applying \eqref{eq:ch21-model-closedness} to this net, with \(\chi=\varphi\)
a model function by hypothesis, yields (1).

Two remarks on Step~6.  It is a genuine net argument: no cofinal sequence in
\(I\) exists once the set of divisorial points is uncountable, and
\Cref{lem:ch21-no-cofinal-sequence} shows that this already happens for
\(\mathbf P^1\) over \(\mathbf C((t))\).  And it uses only clause~(2); the
retraction inequality plays no role here, which is consistent with the fact
that a model function satisfies it automatically.
\end{proof}

\begin{corollary}\label{cor:ch21-reference-independence}
Let \(\|\cdot\|_0\) be a model metric on an ample line bundle \(L\), with
curvature \(\theta\).  A singular metric

\begin{equation}
\label{eq:ch21-singular-metric}
                         \|\cdot\|_\varphi
                         =\|\cdot\|_0\rme^{-\varphi}
\end{equation}

is semipositive precisely when \(\varphi\in\PSH(X,\theta)\), and this
condition is independent of the reference model metric.
\end{corollary}

\begin{proof}
For a model metric, ``semipositive'' means that some determining model line
bundle is relatively nef (\Cref{sec:metrics-semipositive}), and by
\Cref{lem:metrics-nef-independence} this depends only on the metric.  In the
present notation the curvature of \(\|\cdot\|_\varphi\) is
\(\theta+\ddc\varphi\), so the first assertion, for \(\varphi\) a model
function, is exactly \Cref{thm:ch21-model-consistency}.  For general
\(\varphi\), the phrase ``semipositive'' for the singular metric
\eqref{eq:ch21-singular-metric} is \emph{defined} to mean
\(\varphi\in\PSH(X,\theta)\); the content of the corollary is that this
definition does not depend on \(\|\cdot\|_0\).

Suppose then that \(\|\cdot\|_1=\|\cdot\|_0\rme^{-\psi}\) is another model
reference, so that \(\psi\) is a model function.  Its curvature is
\(\theta_1=\theta+\ddc\psi\), while the relative weight of
\(\|\cdot\|_\varphi\) with respect to \(\|\cdot\|_1\) is
\(\varphi-\psi\).  The two total curvature expressions coincide:
\(\theta_1+\ddc(\varphi-\psi)=\theta+\ddc\varphi\).  It remains to see that
\(\varphi\in\PSH(X,\theta)\) if and only if
\(\varphi-\psi\in\PSH(X,\theta_1)\).  Upper semicontinuity and the exclusion
of \(-\infty\) are unaffected, because \(\psi\) is real valued and
continuous.  Clause~(1) is unaffected because \(\psi\) is a model function,
hence, after enlarging \(\mathcal X\) so that \(\psi\) is determined on it,
satisfies \(\psi=\psi\circ p_{\mathcal X}\) by
\Cref{lem:model-function-affine}; subtracting an identity from an inequality
preserves it, and \Cref{prop:ch21-cofinal-independence} allows the
enlargement.  Clause~(2) is unaffected because
\eqref{eq:ch20-reference-change} identifies the \(\theta\)-psh model
functions with \(\psi\) plus the \(\theta_1\)-psh ones, and the bijection
\(u\mapsto u-\psi\) is an isometry of \(C^0(\Delta_{\mathcal X})\) for the
sup-norm; so it carries \(\mathcal P_{\mathcal X}(\theta)\) onto
\(\mathcal P_{\mathcal X}(\theta_1)\) and the uniform limits correspond.
\end{proof}

The corollary is the reason the theory is stated for a form \(\theta\) and a
relative weight \(\varphi\), rather than for a metric: the pair carries
exactly one redundancy, namely the choice of reference, and the corollary
removes it.  The next statement records the two immediate consequences that
are used constantly.

\begin{proposition}
\label{prop:ch21-constants}
Assume \eqref{eq:ch21-hypotheses}.
\begin{enumerate}
\item\label{itm:ch21-const:transl}
If \(\varphi\in\PSH(X,\theta)\) and \(c\in\mathbf R\), then
\(\varphi+c\in\PSH(X,\theta)\).
\item\label{itm:ch21-const:const}
The constant function \(0\) belongs to \(\PSH(X,\theta)\) if and only if
\(\theta\) is semipositive.  In that case \(\PSH(X,\theta)\) contains all
constants and all bounded model psh functions.
\item\label{itm:ch21-const:nonempty}
\(\PSH(X,\theta)\) is nonempty; and if \(\theta\) is semipositive then for
every \(\varepsilon>0\) there is a model function
\(\psi\in\PSH(X,\theta)\) with \(-\varepsilon\leq\psi\leq0\).
\end{enumerate}
\end{proposition}

\begin{proof}
\cref{itm:ch21-const:transl}  Adding a constant changes neither upper
semicontinuity nor the finiteness convention, commutes with
\(p_{\mathcal X}\), and translates the uniform approximants by the same
constant, which are again \(\theta\)-psh model functions because
\(\ddc c=0\) by \eqref{eq:ch20-ddc-constant}.

\cref{itm:ch21-const:const}  The constant \(0\) is a model function, namely
\(\varphi_0\) for the zero divisor.  By
\Cref{thm:ch21-model-consistency} it lies in \(\PSH(X,\theta)\) if and only
if \(\theta+\ddc0=\theta\) is a semipositive model form.  The supplement
follows from \cref{itm:ch21-const:transl} and
\Cref{thm:ch21-model-consistency}.

\cref{itm:ch21-const:nonempty}  Since \(\{\theta\}\) is ample,
\Cref{prop:ch20-positive-representative} provides a model function \(\psi\)
such that \(\theta+\ddc\psi\) is \(\mathcal X\)-positive for a suitable
model, in particular semipositive; by
\Cref{thm:ch21-model-consistency}, \(\psi\in\PSH(X,\theta)\).  If
\(\theta\) is already semipositive, the same proposition provides such a
\(\psi\) with \(-\varepsilon\leq\psi\leq0\), which is
\eqref{eq:ch20-small-positive-representative}.
\end{proof}

\begin{warning}{Ampleness is used, and it is used here}
Part \cref{itm:ch21-const:nonempty} of
\Cref{prop:ch21-constants} is the only place in this chapter where
ampleness of \(\{\theta\}\) is genuinely needed, and it is needed for
nonemptiness alone.  Every other statement of the chapter is a conditional
statement about members of \(\PSH(X,\theta)\) and would remain true, though
possibly vacuous, for a general closed model form.  A reader tracking
hypotheses through \Cref{part:bfj} should therefore expect ampleness to
reappear whenever an \emph{existence} assertion is made --- envelopes,
solutions of a Monge--Amp\`ere equation, finite-energy potentials --- and
not otherwise.
\end{warning}

\section{Equivalent definitions on SNC skeleta}
\label{sec:ch21-equivalent}

\Cref{def:ch21-bfj-psh} is one of several equivalent packagings of the same
class, and different arguments in \Cref{part:bfj} are most easily run from
different ones.  This section collects the packagings and proves their
equivalence, keeping visible which implications are formal and which use the
closedness statement \eqref{eq:ch21-model-closedness} and therefore the
residue characteristic.

\subsection{The equivalence theorem}
\label{sec:ch21-equivalence-theorem}

\begin{theorem}
\label{thm:ch21-equivalent}
Assume \eqref{eq:ch21-hypotheses} and fix a cofinal directed family
\(\mathfrak M\) of SNC models on which \(\theta\) is determined.  Let
\(\varphi:X^{\An}\to[-\infty,+\infty)\) be a function which is not
identically \(-\infty\).  The following are equivalent.

\begin{enumerate}
\item\label{itm:ch21-eq:def}
\(\varphi\in\PSH(X,\theta)\);
\item\label{itm:ch21-eq:cofinal}
\(\varphi\) is upper semicontinuous and the two clauses of
\Cref{def:ch21-bfj-psh} hold for every \(\mathcal X\in\mathfrak M\);
\item\label{itm:ch21-eq:inf}
for every \(\mathcal X\in\mathfrak M\) one has
\(\varphi|_{\Delta_{\mathcal X}}\in\mathcal P_{\mathcal X}(\theta)\), and
\begin{equation}
\label{eq:ch21-inf-of-retractions}
 \varphi=\inf_{\mathcal X\in\mathfrak M}\varphi\circ p_{\mathcal X}
 \qquad\text{pointwise on }X^{\An};
\end{equation}
\item\label{itm:ch21-eq:divisorial}
\(\varphi\) satisfies \eqref{eq:ch21-inf-of-retractions}, the restriction
\(\varphi|_{\Delta_{\mathcal X}}\) is continuous for every
\(\mathcal X\in\mathfrak M\), and there is a net \((\chi_i)\) of
\(\theta\)-psh model functions with \(\chi_i\to\varphi\) pointwise at every
divisorial point;
\item\label{itm:ch21-eq:metric}
\(\varphi\) is upper semicontinuous and the singular metric
\(\|\cdot\|_0\rme^{-\varphi}\) attached to any model metric
\(\|\cdot\|_0\) of curvature \(\theta\) on an ample line bundle is
semipositive in the sense of \Cref{cor:ch21-reference-independence}.
\end{enumerate}
\end{theorem}

\begin{proof}
\cref{itm:ch21-eq:def}\(\Leftrightarrow\)\cref{itm:ch21-eq:cofinal} is
\Cref{prop:ch21-cofinal-independence}, together with the trivial
implication that the clauses for all models imply the clauses for those in
\(\mathfrak M\).

\cref{itm:ch21-eq:def}\(\Rightarrow\)\cref{itm:ch21-eq:inf}.  The membership
\(\varphi|_{\Delta_{\mathcal X}}\in\mathcal P_{\mathcal X}(\theta)\) is
clause~(2), rewritten with \eqref{eq:ch21-skeletal-closure}.  Formula
\eqref{eq:ch21-inf-of-retractions} is
\eqref{eq:ch21-skeletal-reconstruction}, proved in
\Cref{thm:bfj-skeletal-characterization}; the infimum may be taken over the
cofinal family alone because the net \((\varphi\circ p_{\mathcal X})\) is
decreasing and a decreasing net has the same infimum as any cofinal
subnet.

\cref{itm:ch21-eq:inf}\(\Rightarrow\)\cref{itm:ch21-eq:cofinal}.  Assume
\eqref{eq:ch21-inf-of-retractions} and the skeletal membership.  Each
\(\varphi|_{\Delta_{\mathcal X}}\) is continuous, being a uniform limit of
continuous functions on a compact polyhedron, so each
\(\varphi\circ p_{\mathcal X}\) is continuous on \(X^{\An}\).  Hence
\(\varphi\) is an infimum of continuous functions and is upper
semicontinuous by \Cref{lem:usc-inf}.  Clause~(1) holds because
\(\varphi\leq\varphi\circ p_{\mathcal X}\) is one of the terms of the
infimum \eqref{eq:ch21-inf-of-retractions}.  Clause~(2) is the assumed
membership.

\cref{itm:ch21-eq:inf}\(\Rightarrow\)\cref{itm:ch21-eq:divisorial}.
Continuity of the restrictions was just observed.  For the net, run Step~6
of the proof of \Cref{thm:ch21-model-consistency} verbatim: it uses only
clause~(2), which is the assumed membership.

\cref{itm:ch21-eq:divisorial}\(\Rightarrow\)\cref{itm:ch21-eq:inf}.  Fix
\(\mathcal X\in\mathfrak M\).  We must show
\(\varphi|_{\Delta_{\mathcal X}}\in\mathcal P_{\mathcal X}(\theta)\).  Let
\(\Delta=\Delta_{\mathcal X}\) and let \(\Delta^{\divis}\subset\Delta\) be
the set of divisorial points, that is, the points with rational
coordinates; it is dense in \(\Delta\) by
\Cref{rem:quasimonomial-vs-divisorial}\,(1).  By hypothesis
\(\chi_i\to\varphi\) pointwise on \(\Delta^{\divis}\).  By
\Cref{thm:bfj-skeletal-equicontinuity}\,(1) the restrictions
\(\chi_i|_\Delta\) have a common Lipschitz constant \(C\) on each face; and
they are uniformly bounded on \(\Delta\), because they converge at one
point of \(\Delta^{\divis}\) and \(\Delta\) has finite diameter.  A
uniformly Lipschitz family converging pointwise on a dense subset of a
compact metric space converges uniformly on that space: given
\(\varepsilon>0\), cover \(\Delta\) by finitely many balls of radius
\(\varepsilon/(3C)\) centred at points of \(\Delta^{\divis}\), choose
\(i_0\) beyond which all centres are approximated within
\(\varepsilon/3\), and combine.  Hence
\(\chi_i|_\Delta\to\varphi|_\Delta\) uniformly, which is exactly
\(\varphi|_\Delta\in\mathcal P_{\mathcal X}(\theta)\).

\cref{itm:ch21-eq:def}\(\Leftrightarrow\)\cref{itm:ch21-eq:metric} is
\Cref{cor:ch21-reference-independence}.
\end{proof}

\begin{remark}[Which implications are deep]
\label{rem:ch21-which-deep}
Only two steps above are not formal.  The implication
\cref{itm:ch21-eq:divisorial}\(\Rightarrow\)\cref{itm:ch21-eq:inf} uses the
skeletal equicontinuity \Cref{thm:bfj-skeletal-equicontinuity}, proved in
\Cref{chap:bfj-compactness} from an intersection-theoretic estimate on
models; and \Cref{thm:ch21-model-consistency}, used to know that the class
of \(\theta\)-psh model functions is the right one to approximate with, uses
the multiplier ideals of \Cref{sec:ch21-multiplier-input} and therefore
\(\operatorname{char}\widetilde K=0\).  Everything else in
\Cref{thm:ch21-equivalent} is bookkeeping with the retraction identities
\eqref{eq:ch21-retraction-composition} and the elementary topology of
\Cref{chap:topological-tools}.  In particular the equivalence of
\cref{itm:ch21-eq:def} and \cref{itm:ch21-eq:inf} holds verbatim over any
complete discretely valued field for which SNC models are cofinal.
\end{remark}

\subsection{The criterion on a curve, made finite}
\label{sec:ch21-equivalent-curve}

In dimension one the skeletal condition becomes a finite list of numerical
inequalities at the vertices of a graph, and it is worth recording that
translation once, both because every worked example of
\Cref{sec:ch21-worked} uses it and because it displays the exact point at
which the higher-dimensional theory ceases to be finite.

\begin{proposition}
\label{prop:ch21-curve-criterion}
Let \(K\) be complete and discretely valued, let \(C\) be a smooth
projective connected \(K\)-curve, and let \(\theta\) be a closed model form
determined on a semistable model \(\mathcal C\), with dual complex
\(\Gamma\) and curvature measure \(\mu_0\) as in
\eqref{eq:graph-curvature}.  Let
\(\varphi:C^{\An}\to[-\infty,+\infty)\) be a function, not identically
\(-\infty\).  Then \(\varphi\in\PSH(C,\theta)\) if and only if
\begin{enumerate}
\item\label{itm:ch21-cc:inf}
\(\varphi=\inf_{\mathcal C'}\varphi\circ p_{\mathcal C'}\), the infimum
being over the SNC models \(\mathcal C'\geq\mathcal C\) --- not merely the
semistable ones, which are not cofinal
(\Cref{rem:ch21-semistable-not-cofinal}); and
\item\label{itm:ch21-cc:pl}
for every such \(\mathcal C'\), with skeleton \(\Gamma'\) carrying its
canonical metric (\Cref{rem:ch20-multiplicities-lengths}), the restriction
\(\varphi|_{\Gamma'}\) is a uniform limit of piecewise affine functions
\(u\) on the metric graph \(\Gamma'\) satisfying the finitely many
inequalities
\begin{equation}
\label{eq:ch21-curve-vertex-inequalities}
 \mu_0(\{q\})+\sum_{\nu\in T_q\Gamma'}d_\nu u(q)\ \geq\ 0
 \qquad\text{for every }q\in\Gamma' .
\end{equation}
\end{enumerate}
Here \(d_\nu u(q)\) is the outgoing slope \eqref{eq:ch08-outgoing-slope} and
only the finitely many vertices and slope-break points of \(u\) give a
nontrivial condition.
\end{proposition}

\begin{proof}
By \Cref{thm:ch21-equivalent},
\cref{itm:ch21-eq:def}\(\Leftrightarrow\)\cref{itm:ch21-eq:inf}, membership is
equivalent to \cref{itm:ch21-cc:inf} together with
\(\varphi|_{\Gamma'}\in\mathcal P_{\mathcal C'}(\theta)\) for every
\(\mathcal C'\) in the cofinal family of SNC models dominating
\(\mathcal C\), which is cofinal by \Cref{prop:snc-cofinal}.  It therefore
suffices to identify
\(\mathcal P_{\mathcal C'}(\theta)\) with the uniform closure of the
piecewise affine functions satisfying
\eqref{eq:ch21-curve-vertex-inequalities}.

A \(\theta\)-psh model function \(\psi\), determined on some
\(\mathcal C''\geq\mathcal C'\) with skeleton
\(\Gamma''\supseteq\Gamma'\), restricts to a piecewise affine function on
\(\Gamma'\) by \Cref{prop:skeleton-model-pa}.  On \(\Gamma''\) itself,
\eqref{eq:ch20-curve-laplacian-general} identifies the relative degree of
the associated model line bundle on a component \(E\) with
\(b_E^{-1}(\Delta_{\Gamma''}\psi)(\{x_E\})\), the factor \(b_E>0\) being
irrelevant to the sign (when \(\mathcal C''\) is semistable this is
\Cref{prop:ch08-curve-cl} with \(b_E=1\)); together with
\eqref{eq:graph-curvature} and the nef test
\eqref{eq:ch20-curve-nef-test}, relative nefness is exactly
\(\mu_0+\Delta_{\Gamma''}\psi\geq0\) on \(\Gamma''\).

That identity does \emph{not} pass to the restriction: the model functions
occurring in \(\mathcal P_{\mathcal C'}(\theta)\) are not required to be
determined on \(\mathcal C'\), and for one determined on a higher model the
Laplacian of \(\psi|_{\Gamma'}\) is in general strictly larger than the
vector of relative degrees --- \Cref{ex:ch21-conic-refinement} is the
smallest instance.  What passes is the inequality.  Fix \(q\in\Gamma'\) and
decompose \(\Gamma''\setminus\Gamma'\) into connected components, each
attached to \(\Gamma'\) at a single point; then
\[
 \bigl(\Delta_{\Gamma''}\psi\bigr)(\{q\})
 =\bigl(\Delta_{\Gamma'}(\psi|_{\Gamma'})\bigr)(\{q\})
  +\sum_{T\ni q}\ \sum_{e\subset T,\ e\ni q}\partial_e\psi(q),
\]
the outer sum running over the components \(T\) attached at \(q\).  Each
inner sum is \(\leq0\) by the flux inequality
\eqref{eq:ch21-curve-flux}, proved in Step~2 of
\Cref{prop:ch21-curve-retraction} below from
\eqref{eq:ch20-nef-pullback-test} and
\eqref{eq:graph-laplacian} alone, hence without reference to the present
proposition.  Since \(\mu_0\) gains no mass on a higher skeleton
\eqref{eq:ch21-curve-no-new-mass}, we conclude
\[
 \bigl(\Delta_{\Gamma'}(\psi|_{\Gamma'})\bigr)(\{q\})
 \ \geq\ \bigl(\Delta_{\Gamma''}\psi\bigr)(\{q\})
 \ \geq\ -\mu_0(\{q\}),
\]
which is \eqref{eq:ch21-curve-vertex-inequalities}.  Conversely, a piecewise affine
\(u\) on \(\Gamma'\) satisfying
\eqref{eq:ch21-curve-vertex-inequalities} is, after a rational
approximation of its breakpoints and slopes, the restriction of a model
function by \Cref{thm:skeleton-pl-realization}, and that model function is
\(\theta\)-psh by the same identification; the rational approximation may be
taken uniformly small and preserves
\eqref{eq:ch21-curve-vertex-inequalities} up to an arbitrarily small error,
which is absorbed by the uniform closure.  Taking uniform closures on both
sides completes the identification.
\end{proof}

\begin{remark}[Semistable models are not cofinal]
\label{rem:ch21-semistable-not-cofinal}
Over a fixed complete discretely valued \(K\) the semistable models of
\(C\) do \emph{not} form a cofinal family, and no blowup repairs this: a
blowup can only raise multiplicities, never lower them.  The running conic
is already a witness.  Blowing up the node of \(\mathcal C=\{xy=\pi\}\)
produces a central component \(E_{1/2}\) of multiplicity \(2\)
(\Cref{ex:ch21-conic-numbers}), and blowing up a closed point of
\(E_{1/2}\) away from the nodes produces a further component \(E'\), again
of multiplicity \(2\), whose divisorial point \(x_{E'}\) lies on no
semistable skeleton: a point of \(C^{\An}\) is a \emph{vertex} of a semistable skeleton only if it
is divisorial with multiplicity \(1\), and every point of an \emph{edge} lies
on the geodesic joining two such vertices.  Multiplicity alone is therefore
no obstruction: \(x_{E_{1/2}}\) itself has multiplicity \(2\) and is the
midpoint of the semistable skeleton \(\Gamma_0\).  What excludes \(x_{E'}\)
is its centre.  In the case at hand, \(C\) a conic over \(K=k((t))\) with
\(k\) algebraically closed, the multiplicity-one divisorial points are the
Gauss points of \(K\)-rational discs, so every point of a semistable edge is
of the form \(\zeta_{c,r}\) with \(c\in K\); whereas \(x_{E'}\) is the Galois
orbit of a pair of discs defined only over the ramified quadratic extension
\(K(\pi^{1/2})\).  Removing the
obstruction requires a ramified base extension \(K'/K\), which is the
content of the semistable reduction theorem and is exactly what is
\emph{not} available here, since \(\theta\), \(L\) and \(\PSH(C,\theta)\)
are attached to \(K\).  Accordingly
\Cref{prop:ch21-curve-criterion} quantifies over SNC models, which
\emph{are} cofinal (\Cref{prop:snc-cofinal}), at the price of the
multiplicities in \eqref{eq:ch20-curve-laplacian-general}.  Semistable
models remain the right family for \emph{computations}---all of
\Cref{sec:ch21-worked} stays inside them, and there the edge lengths are
\(1\) and the weights disappear---but they are not the right family for a
quantifier.
\end{remark}

\begin{remark}[What is finite and what is not]
\label{rem:ch21-curve-finite}
For a \emph{fixed} skeleton and a fixed piecewise affine \(u\),
\eqref{eq:ch21-curve-vertex-inequalities} is a finite system of linear
inequalities, one per vertex, and it may be checked by hand; every
computation in \Cref{sec:ch21-worked} is such a check.  What is not finite is
the quantifier over models: \(\Gamma'\) ranges over an uncountable directed
family (\Cref{lem:ch21-no-cofinal-sequence}), and clause
\cref{itm:ch21-cc:inf} is a statement about all of them at once.  In
dimension \(n\geq2\) even the first finiteness fails: relative nefness is
tested against the vertical curves of an \(n\)-dimensional special fibre,
which are not indexed by the vertices of the dual complex, and no finite
vertex criterion is available.  This is the same phenomenon as item (L4) of
\Cref{sec:ch19-limitations-table}, seen from the model side.
\end{remark}

\subsection{Upper semicontinuity as a skeletal statement}
\label{sec:ch21-usc-skeletal}

The equivalence \cref{itm:ch21-eq:cofinal}\(\Leftrightarrow\)
\cref{itm:ch21-eq:inf} of \Cref{thm:ch21-equivalent} has an instructive
corollary: once the skeletal approximation clause is granted, upper
semicontinuity and the reconstruction formula are the same condition.  This
is worth isolating, because it is the mechanism by which every
\(-\infty\)-value in the theory is produced.

\begin{proposition}
\label{prop:ch21-usc-iff-inf}
Let \(\varphi:X^{\An}\to[-\infty,+\infty)\) be a function, not identically
\(-\infty\), satisfying clause~(2) of \Cref{def:ch21-bfj-psh} for every SNC
model determining \(\theta\), and satisfying the retraction inequality
\eqref{eq:ch21-retraction-inequality}.  Then
\[
 \varphi\ \text{is upper semicontinuous}
 \quad\Longleftrightarrow\quad
 \varphi=\inf_{\mathcal X}\varphi\circ p_{\mathcal X}.
\]
\end{proposition}

\begin{proof}
\((\Leftarrow)\)  Each \(\varphi|_{\Delta_{\mathcal X}}\) is continuous by
clause~(2) and each \(p_{\mathcal X}\) is continuous by
\Cref{prop:skeleton-retraction}, so \(\varphi\circ p_{\mathcal X}\in
C^0(X^{\An})\).  An infimum of continuous functions is upper
semicontinuous (\Cref{lem:usc-inf}).

\((\Rightarrow)\)  The retraction inequality gives
\(\varphi\leq\inf_{\mathcal X}\varphi\circ p_{\mathcal X}=:h\).  For the
converse inequality fix \(x\in X^{\An}\).  By
\Cref{thm:skeleton-inverse-limit} the net \((p_{\mathcal X}(x))_{\mathcal X}\)
converges to \(x\), so by the net test for upper semicontinuity
\Cref{prop:net-semicontinuity},
\[
 \varphi(x)\ \geq\ \varlimsup_{\mathcal X}\varphi(p_{\mathcal X}(x)).
\]
The net \((\varphi(p_{\mathcal X}(x)))_{\mathcal X}\) is decreasing by
\eqref{eq:ch21-decreasing-retractions}, hence converges to its infimum
\(h(x)\) by \Cref{lem:monotone-net-converges}, and its upper limit is
\(h(x)\).  Therefore \(\varphi(x)\geq h(x)\).
\end{proof}

\begin{remark}[What the proposition does not say]
\label{rem:ch21-usc-not-continuity}
It does not say that \(\varphi\) is continuous.  The functions
\(\varphi\circ p_{\mathcal X}\) are continuous and decrease to \(\varphi\),
but a decreasing net of continuous functions on a compact space converges
uniformly only when its limit is continuous, by
\Cref{prop:dini-nets}; and the limit here is typically not continuous, as
\Cref{ex:ch21-conic-refinement} shows with a function that decreases to
\(-\infty\) at two rigid points.  The correct reading is the
contrapositive: \(\varphi\) is continuous if and only if the decreasing net
of its skeletal shadows converges uniformly, and that is exactly the
situation of \Cref{thm:ch21-continuous-zhang}.
\end{remark}

\section{Restriction to dual complexes and skeletal reconstruction}
\label{sec:ch21-restriction}

The compatibility \eqref{eq:ch21-retraction-composition} of the retractions
is what makes the family \((\varphi\circ p_{\mathcal X})_{\mathcal X}\)
monotone, and monotonicity is what converts the inverse limit
\Cref{thm:skeleton-inverse-limit} into a reconstruction formula.  We begin
with the reconstruction theorem itself and then examine the resulting
inverse system of restrictions.

\subsection{The reconstruction theorem}
\label{sec:ch21-reconstruction}

\begin{theorem}\label{thm:bfj-skeletal-characterization}
If \(\varphi\in\PSH(X,\theta)\), then the following statements hold.

\begin{enumerate}
\item For every determining SNC model, the restriction
\(\varphi|_{\Delta_{\mathcal X}}\) is finite, continuous, and convex on
each face.
\item As the model increases, the continuous functions
\(\varphi\circ p_{\mathcal X}\) form a decreasing net and

\begin{equation}
\label{eq:ch21-skeletal-reconstruction}
              \varphi
              =\inf_{\mathcal X}\varphi\circ p_{\mathcal X}
              =\lim_{\mathcal X}\varphi\circ p_{\mathcal X}
              \quad\text{pointwise on }X^{\An}.
\end{equation}

\item The function is uniquely determined by its values at divisorial
points.
\end{enumerate}

Conversely, the two clauses of \Cref{def:ch21-bfj-psh}, checked on a
cofinal SNC system, characterize the class.  Convexity on faces by itself is
not a converse.
\end{theorem}

\begin{proof}
\emph{Step 1: the restrictions.}
The first assertion was observed after the definition and proved as
\Cref{prop:ch21-first-consequences}\,\cref{itm:ch21-fc:finite}.

\emph{Step 2: the retractions decrease.}
Suppose
\(\mathcal X'\geq\mathcal X\).  Applying
\eqref{eq:ch21-retraction-inequality} at the point
\(p_{\mathcal X'}(x)\) and using
\eqref{eq:ch21-retraction-composition} gives

\begin{equation}
\label{eq:ch21-decreasing-retractions}
 \varphi(x)
 \leq\varphi(p_{\mathcal X'}(x))
 \leq\varphi(p_{\mathcal X}(x)).
\end{equation}

The two inequalities in \eqref{eq:ch21-decreasing-retractions} come from
different sources: the first is clause~(1) for \(\mathcal X'\); the second is
clause~(1) for \(\mathcal X\), applied at the point
\(p_{\mathcal X'}(x)\), combined with
\(p_{\mathcal X}p_{\mathcal X'}=p_{\mathcal X}\).  The family
\((\varphi\circ p_{\mathcal X})_{\mathcal X}\) is therefore a decreasing net
of continuous functions indexed by the directed set of determining SNC
models.

\emph{Step 3: the reconstruction formula.}
Let the middle term decrease to \(\widetilde\varphi(x)\); it exists by
\Cref{lem:monotone-net-converges} and equals
\(\inf_{\mathcal X}\varphi(p_{\mathcal X}(x))\).  It is at least
\(\varphi(x)\).  On the other hand, \(p_{\mathcal X}(x)\to x\) as
the model increases, by \Cref{thm:skeleton-inverse-limit}, and upper
semicontinuity gives, through the net test
\Cref{prop:net-semicontinuity},

\[
 \varphi(x)\geq
 \varlimsup_{\mathcal X}\varphi(p_{\mathcal X}(x))
 =\widetilde\varphi(x).
\]

This proves \eqref{eq:ch21-skeletal-reconstruction}.  Observe that both
displayed limits are limits of \emph{nets}: the index set is the set of
determining SNC models, and \Cref{lem:ch21-no-cofinal-sequence} forbids
replacing it by a sequence.

\emph{Step 4: divisorial determination.}
Divisorial points are exactly the points of a dual complex with rational
coordinates (\Cref{rem:quasimonomial-vs-divisorial}\,(1)), hence they are
dense in every dual complex.
If two psh functions agree on all divisorial points, continuity of their
restrictions, which is Step~1, gives
equality on each skeleton; then
\eqref{eq:ch21-skeletal-reconstruction} gives equality globally.  The final
assertion is \Cref{prop:ch21-cofinal-independence}; that convexity on faces
is not a converse is \Cref{ex:ch21-clause-one-needed}, whose function is
convex, indeed constant, on every face.
\end{proof}

\begin{proposition}\label{prop:ch21-divisorial-order}
If \(\varphi,\psi\in\PSH(X,\theta)\) and
\(\varphi\leq\psi\) at every divisorial point, then
\(\varphi\leq\psi\) on \(X^{\An}\).
\end{proposition}

\begin{proof}
Density of the divisorial points in each dual complex and continuity of the
two restrictions, by
\Cref{prop:ch21-first-consequences}\,\cref{itm:ch21-fc:finite}, imply the
inequality on every dual complex.  Hence,
for every \(x\),

\[
 \varphi(x)
 =\inf_{\mathcal X}\varphi(p_{\mathcal X}(x))
 \leq\inf_{\mathcal X}\psi(p_{\mathcal X}(x))
 =\psi(x),
\]

by \eqref{eq:ch21-skeletal-reconstruction}.
\end{proof}

The same argument also locates the supremum.

\begin{proposition}\label{prop:ch21-sup-on-skeleton}
For every \(\varphi\in\PSH(X,\theta)\) and every SNC model on which
\(\theta\) is determined,

\begin{equation}
\label{eq:ch21-sup-on-skeleton}
             \sup_{X^{\An}}\varphi
             =\max_{\Delta_{\mathcal X}}\varphi.
\end{equation}
\end{proposition}

\begin{proof}
Inequality \eqref{eq:ch21-retraction-inequality} bounds every value by a
skeletal value.  The reverse inequality follows from
\(\Delta_{\mathcal X}\subset X^{\An}\).  Continuity and compactness of the
dual complex turn its supremum into a maximum.
\end{proof}

\begin{remark}[A sharper statement than compactness gives]
\label{rem:ch21-sup-sharper}
Upper semicontinuity on the compact space \(X^{\An}\) already gives
\(\sup\varphi=\max\varphi\) and the existence of a maximizer
(\Cref{cor:usc-attains-max}).  The content of
\Cref{prop:ch21-sup-on-skeleton} is that a maximizer may be found on
\emph{any} prescribed determining dual complex, a finite polyhedron of
dimension at most \(n\).  This is what makes normalizations such as
\(\max_{\Delta_{\mathcal X}}\varphi=0\), used throughout
\Cref{chap:bfj-compactness}, computable, and it is the first place where the
retraction inequality does work that compactness cannot do.
\end{remark}

\subsection{The inverse system of restrictions}
\label{sec:ch21-inverse-system}

\Cref{thm:bfj-skeletal-characterization} says that a psh function is
determined by the family of its restrictions.  It is useful to record the
structure of that family precisely, since \Cref{chap:bfj-compactness}
topologizes \(\PSH(X,\theta)\) through it.

\begin{proposition}
\label{prop:ch21-restriction-system}
For SNC models \(\mathcal X'\geq\mathcal X\) determining \(\theta\), write
\(\pi_{\mathcal X'\mathcal X}:=p_{\mathcal X}|_{\Delta_{\mathcal X'}}:
\Delta_{\mathcal X'}\to\Delta_{\mathcal X}\).  Then:
\begin{enumerate}
\item\label{itm:ch21-rs:transition}
\(\pi_{\mathcal X'\mathcal X}\) is a continuous surjection restricting to
the identity on \(\Delta_{\mathcal X}\subset\Delta_{\mathcal X'}\), and
\(\pi_{\mathcal X''\mathcal X}=\pi_{\mathcal X'\mathcal X}\circ
\pi_{\mathcal X''\mathcal X'}\) for
\(\mathcal X''\geq\mathcal X'\geq\mathcal X\);
\item\label{itm:ch21-rs:compat}
for \(\varphi\in\PSH(X,\theta)\), the restrictions satisfy
\begin{equation}
\label{eq:ch21-restriction-compatibility}
 \varphi|_{\Delta_{\mathcal X}}
 =\bigl(\varphi|_{\Delta_{\mathcal X'}}\bigr)\big|_{\Delta_{\mathcal X}},
 \qquad
 \varphi|_{\Delta_{\mathcal X'}}
 \leq\bigl(\varphi|_{\Delta_{\mathcal X}}\bigr)\circ
     \pi_{\mathcal X'\mathcal X};
\end{equation}
\item\label{itm:ch21-rs:injective}
the map
\(\varphi\mapsto(\varphi|_{\Delta_{\mathcal X}})_{\mathcal X}\) from
\(\PSH(X,\theta)\) to \(\prod_{\mathcal X}C^0(\Delta_{\mathcal X})\) is
injective, with left inverse given by
\eqref{eq:ch21-skeletal-reconstruction};
\item\label{itm:ch21-rs:notsurjective}
the map of \cref{itm:ch21-rs:injective} is not surjective onto the compatible
systems: a family \((u_{\mathcal X})\) satisfying
\eqref{eq:ch21-restriction-compatibility} with
\(u_{\mathcal X}\in\mathcal P_{\mathcal X}(\theta)\) determines at most one
candidate \(\varphi:=\inf_{\mathcal X}u_{\mathcal X}\circ p_{\mathcal X}\),
and this candidate lies in \(\PSH(X,\theta)\) if and only if it restricts
back to \(u_{\mathcal X}\) on each \(\Delta_{\mathcal X}\).
\end{enumerate}
\end{proposition}

\begin{proof}
\cref{itm:ch21-rs:transition}  Continuity is
\Cref{prop:skeleton-retraction}; that \(p_{\mathcal X}\) is the identity on
\(\Delta_{\mathcal X}\) is the same proposition, and surjectivity follows
because \(\Delta_{\mathcal X}\subset\Delta_{\mathcal X'}\) is fixed
pointwise.  The cocycle identity is
\(p_{\mathcal X}p_{\mathcal X'}=p_{\mathcal X}\), which is
\eqref{eq:ch21-retraction-composition}.

\cref{itm:ch21-rs:compat}  The first identity is the inclusion
\(\Delta_{\mathcal X}\subset\Delta_{\mathcal X'}\).  The inequality is
clause~(1) for \(\mathcal X\) restricted to \(\Delta_{\mathcal X'}\).

\cref{itm:ch21-rs:injective}  This is
\Cref{thm:bfj-skeletal-characterization}\,(3) together with
\eqref{eq:ch21-skeletal-reconstruction}.

\cref{itm:ch21-rs:notsurjective}  Uniqueness of the candidate is
\eqref{eq:ch21-skeletal-reconstruction} again.  Suppose
\(\varphi:=\inf_{\mathcal X}u_{\mathcal X}\circ p_{\mathcal X}\) restricts
to \(u_{\mathcal X}\) on each \(\Delta_{\mathcal X}\).  Then \(\varphi\)
satisfies \eqref{eq:ch21-inf-of-retractions} by construction and
\(\varphi|_{\Delta_{\mathcal X}}=u_{\mathcal X}\in
\mathcal P_{\mathcal X}(\theta)\) by hypothesis; and \(\varphi\) is not
identically \(-\infty\), because on one fixed \(\Delta_{\mathcal X_0}\) it
agrees with the real-valued continuous function \(u_{\mathcal X_0}\).
Hence \Cref{thm:ch21-equivalent}\,\cref{itm:ch21-eq:inf} applies and
\(\varphi\in\PSH(X,\theta)\).  The converse implication is
\eqref{eq:ch21-skeletal-reconstruction}.
\end{proof}

\begin{warning}{A compatible system need not come from a psh function}
The last part of \Cref{prop:ch21-restriction-system} is a genuine
restriction and not a formality.  The consistency conditions
\eqref{eq:ch21-restriction-compatibility} only constrain a system along
inclusions of skeleta; nothing in them prevents the candidate
\(\inf_{\mathcal X}u_{\mathcal X}\circ p_{\mathcal X}\) from dropping
strictly below \(u_{\mathcal X_0}\) at a point of
\(\Delta_{\mathcal X_0}\), because the infimum is taken over models whose
skeleta pass through that point along new branches.  Concretely: on the
curve of \Cref{ex:ch21-conic-refinement}, prescribing \(u\equiv0\) on
\(\Delta_{\mathcal C}\) and \(u'(\tau)=-\tau\) on the hair of
\(\Delta_{\mathcal C_1}\) is consistent in the sense of
\eqref{eq:ch21-restriction-compatibility}, and it does extend; but
prescribing \(u'(\tau)=+\tau\) on the hair violates the second inequality of
\eqref{eq:ch21-restriction-compatibility} and no psh function realizes it.
The moral is that the system of restrictions must be checked against
clause~(1) at every stage, not only against convexity.
\end{warning}

\subsection{Restriction of the retraction: the curve case}
\label{sec:ch21-curve-retraction}

There is one natural operation that the reader will be tempted to perform
and that the book does \emph{not} assert in general: replacing
\(\varphi\) by \(\varphi\circ p_{\mathcal X}\) and expecting the result to
be \(\theta\)-psh.  It is continuous and it decreases to \(\varphi\), so if
it were always psh, the decreasing regularization theorem
\Cref{thm:bfj-decreasing-regularization} would reduce at once to its
continuous case.  It is therefore worth knowing exactly when the operation
is legitimate.  On curves it always is, and the proof is a two-line flux
computation which we give in full because it is the model for every later
skeletal estimate.

\begin{proposition}
\label{prop:ch21-curve-retraction}
Let \(K\) be complete and discretely valued, let \(C\) be a smooth
projective connected \(K\)-curve, let \(\theta\) be a closed model form
determined on a semistable model \(\mathcal C\) with dual complex
\(\Gamma=\Delta_{\mathcal C}\) and retraction \(p=p_{\mathcal C}\), and let
\(\mu_0\) be the curvature measure of \(\theta\) on \(\Gamma\) in the sense
of \eqref{eq:graph-curvature}.  Recall from
\eqref{eq:ch20-curve-laplacian-general} and \eqref{eq:ch20-curve-nef-test}
that a model
function \(\psi\) determined on an SNC model with skeleton
\(\Gamma'\supseteq\Gamma\), the latter carrying its canonical metric
(\Cref{rem:ch20-multiplicities-lengths}), is \(\theta\)-psh if and only if
\begin{equation}
\label{eq:ch21-curve-psh-test}
 \mu_0+\Delta_{\Gamma'}\psi\ \geq\ 0
 \qquad\text{as a measure on }\Gamma' ,
\end{equation}
with \(\Delta_{\Gamma'}\) the graph Laplacian
\eqref{eq:graph-laplacian}; here \(\mu_0\) is regarded as a measure on
\(\Gamma'\) through \(\Gamma\subseteq\Gamma'\), and the multiplicities
\(b_E\) of \eqref{eq:ch20-curve-laplacian-general}, being positive, do not
affect the inequality (for \(\mathcal C'\) semistable they are all \(1\) and
the criterion is \Cref{prop:ch08-curve-cl}).  Then, for every
\(\varphi\in\PSH(C,\theta)\),
\[
 \varphi\circ p\ \in\ \PSH(C,\theta)\cap C^0(C^{\An}),
\]
and consequently every \(\theta\)-psh function on a curve is the decreasing
limit of a net of \emph{continuous} \(\theta\)-psh functions.
\end{proposition}

\begin{proof}
\emph{Step 1: the reference measure does not charge the new part.}  Let
\(\mathcal C'\geq\mathcal C\) be an SNC model with skeleton
\(\Gamma'\supseteq\Gamma\) and let \(\rho:\mathcal C'\to\mathcal C\) be the
domination.  Since \(\theta\) is determined on \(\mathcal C\), its
determination on \(\mathcal C'\) is \(\rho^*\theta_{\mathcal C}\), whose
degree on a component \(E\) contracted by \(\rho\) is \(0\) by the first
case of \eqref{eq:ch20-nef-pullback-test}.  As the reference measure of a
model form on a curve assigns to the vertex \(x_E\) the mass
\(b_E\deg(\theta_{\mathcal C'}|_E)\), we conclude
\begin{equation}
\label{eq:ch21-curve-no-new-mass}
 \mu_0\bigl(\Gamma'\setminus\Gamma\bigr)=0 .
\end{equation}

\emph{Step 2: the flux inequality.}  Write \(\Gamma'\setminus\Gamma\) as a
disjoint union of connected components; the closure \(T\) of such a
component meets \(\Gamma\) in a single point \(q\), because \(\Gamma'\) is a
graph containing \(\Gamma\) and each component of the complement of a
connected subgraph is attached along one point.  Let \(u\) be piecewise
affine on \(\Gamma'\) with \(\mu_0+\Delta_{\Gamma'}u\geq0\).  Summing
\eqref{eq:graph-laplacian} over the points of \(T\setminus\{q\}\), all
interior edges contribute twice with opposite signs, and one is left with
the net outward flux at \(q\):
\[
 \Delta_{\Gamma'}u\bigl(T\setminus\{q\}\bigr)
 =-\sum_{e\subset T,\ e\ni q}\partial_e u(q),
\]
where \(\partial_e u(q)\) is the outgoing slope at \(q\) along the edge
\(e\).  Combining with \eqref{eq:ch21-curve-no-new-mass} and positivity,
\begin{equation}
\label{eq:ch21-curve-flux}
 \sum_{e\subset T,\ e\ni q}\partial_e u(q)
 =-\Delta_{\Gamma'}u\bigl(T\setminus\{q\}\bigr)
 \leq\mu_0\bigl(T\setminus\{q\}\bigr)=0 .
\end{equation}

\emph{Step 3: the retraction of a model psh function.}  Let \(\psi\) be a
\(\theta\)-psh model function and let \(u=\psi|_{\Gamma}\).  We claim that
\(\psi\circ p\) is again a \(\theta\)-psh model function.  It is a model
function: \(u\) is piecewise affine on \(\Gamma\)
(\Cref{prop:skeleton-model-pa}), so by \Cref{thm:skeleton-subdivision} there
is a vertical blowup of \(\mathcal C\) with the same skeleton \(\Gamma\),
subdivided so that \(u\) is affine on each cell, and \(u\) is then the model
function of a vertical \(\mathbf Q\)-Cartier divisor on it; that model
function equals \(u\circ p=\psi\circ p\) by
\Cref{lem:model-function-affine}.  For positivity, test
\eqref{eq:ch21-curve-psh-test} on an arbitrary skeleton
\(\Gamma'\supseteq\Gamma\).  On \(\Gamma'\setminus\Gamma\) the function
\(\psi\circ p\) is locally constant, so
\(\Delta_{\Gamma'}(\psi\circ p)=0\) there and \(\mu_0\geq0\) suffices.  At a
point \(q\in\Gamma\), the outgoing slopes of \(\psi\circ p\) along edges of
\(\Gamma'\setminus\Gamma\) vanish, while along edges inside \(\Gamma\) they
agree with those of \(u\); hence
\[
 \bigl(\mu_0+\Delta_{\Gamma'}(\psi\circ p)\bigr)(\{q\})
 =\mu_0(\{q\})+\Delta_{\Gamma}u(\{q\})
 \geq\mu_0(\{q\})+\Delta_{\Gamma'}\psi(\{q\})
 \geq0,
\]
the middle inequality being \eqref{eq:ch21-curve-flux} applied to \(\psi\)
on each component of \(\Gamma'\setminus\Gamma\) attached at \(q\).  This
proves the claim.

\emph{Step 4: passage to \(\varphi\).}  Let \(\psi_k\) be \(\theta\)-psh
model functions with \(\psi_k\to\varphi\) uniformly on \(\Gamma\).  Then
\(\psi_k\circ p\to\varphi\circ p\) uniformly on \(C^{\An}\), because
\(p\) has image \(\Gamma\).  By Step~3 each \(\psi_k\circ p\) is a
\(\theta\)-psh model function, so \(\varphi\circ p\) is a uniform limit of
such on all of \(C^{\An}\); in particular its restriction to every dual
complex lies in the corresponding \(\mathcal P(\theta)\), and clause~(1)
holds for \(\varphi\circ p\) since each \(\psi_k\circ p\) satisfies it
(\Cref{lem:ch21-model-skeletal-properties}) and uniform limits preserve
inequalities.  Continuity is clear.  Hence
\(\varphi\circ p\in\PSH(C,\theta)\cap C^0(C^{\An})\).

Finally, \eqref{eq:ch21-skeletal-reconstruction} exhibits \(\varphi\) as the
decreasing limit of the net \((\varphi\circ p_{\mathcal C'})\), each member
of which is continuous and psh by what we have just proved applied to
\(\mathcal C'\).
\end{proof}

\begin{warning}{Not asserted in dimension at least two}
\Cref{prop:ch21-curve-retraction} is a theorem about curves and is used
only in \Cref{sec:ch21-worked}.  For \(\dim X\geq2\) this book makes no
claim that \(\varphi\circ p_{\mathcal X}\) is \(\theta\)-psh, and the reader
should not supply one: the flux computation of Step~2 replaces a global nef
test by finitely many numerical inequalities at vertices, and no such finite
list exists in higher dimension, exactly as in item (L4) of the limitations
table \Cref{sec:ch19-limitations-table}.  The decreasing regularization
theorem \Cref{thm:bfj-decreasing-regularization} is proved by a different
route, through envelopes; the curve case above is not a shortcut to it, but
it does explain why every statement in \Cref{part:bfj} is easier on curves.
\end{warning}

\section{Stability}
\label{sec:ch21-stability}

This section proves the closure properties that make
\(\PSH(X,\theta)\) usable: constants, maxima, convex combinations, sums
across two classes, uniform limits, decreasing nets, and pullback.  Each is
proved separately, and each proof is organized in the same way, because the
same three verifications recur: upper semicontinuity, the retraction
inequality, and clause~(2) on a fixed dual complex.  The last of these is
always the only one with content, and it always reduces to the corresponding
statement for model functions from \Cref{chap:bfj-model-forms} together with
a Lipschitz estimate for the operation in question.

\subsection{The elementary calculus}
\label{sec:ch21-elementary-calculus}

\begin{proposition}\label{prop:ch21-elementary-stability}
The class \(\PSH(X,\theta)\) is convex and is stable under addition of real
constants and finite maxima.
\end{proposition}

\begin{proof}
This is the special case \(\theta_1=\theta_2=\theta\),
\(t_1+t_2=1\), of \Cref{thm:ch21-stability}, together with
\Cref{prop:ch21-constants}\,\cref{itm:ch21-const:transl}.
\end{proof}

The next theorem is the same statement in the generality actually used
later, where the two potentials may be attached to different forms.  We
state and prove it in full because the bookkeeping with \(-\infty\) is where
mistakes occur.

\begin{theorem}
\label{thm:ch21-stability}
Assume \eqref{eq:ch21-hypotheses} for two closed model forms
\(\theta_1,\theta_2\).
\begin{enumerate}
\item\label{itm:ch21-stab:const}
If \(\varphi\in\PSH(X,\theta_1)\) and \(c\in\mathbf R\), then
\(\varphi+c\in\PSH(X,\theta_1)\).
\item\label{itm:ch21-stab:sum}
If \(\varphi_j\in\PSH(X,\theta_j)\) and \(t_1,t_2\geq0\), then
\[
 t_1\varphi_1+t_2\varphi_2\in\PSH(X,t_1\theta_1+t_2\theta_2),
\]
provided \(t_1+t_2>0\).  In particular \(\PSH(X,\theta)\) is convex, and
\(\PSH(X,\theta_1)+\PSH(X,\theta_2)\subseteq\PSH(X,\theta_1+\theta_2)\).
\item\label{itm:ch21-stab:max}
If \(\varphi,\psi\in\PSH(X,\theta_1)\), then
\(\max\{\varphi,\psi\}\in\PSH(X,\theta_1)\); more generally a maximum of
finitely many members of \(\PSH(X,\theta_1)\) is again one.
\item\label{itm:ch21-stab:unif}
If \(\varphi_k\in\PSH(X,\theta_1)\) converge uniformly on \(X^{\An}\) to
\(\varphi\), then \(\varphi\in\PSH(X,\theta_1)\).
\item\label{itm:ch21-stab:mono}
If \(\varphi\in\PSH(X,\theta_1)\) and \(\theta_2-\theta_1\) is semipositive,
then \(\varphi\in\PSH(X,\theta_2)\).
\end{enumerate}
\end{theorem}

\begin{proof}
Throughout, fix an SNC model \(\mathcal X\) determining both
\(\theta_1\) and \(\theta_2\); such a model exists because two models are
dominated by a third, and \Cref{prop:ch21-cofinal-independence} lets us test
the definition on the cofinal family of models dominating it.

\cref{itm:ch21-stab:const} is
\Cref{prop:ch21-constants}\,\cref{itm:ch21-const:transl}.

\cref{itm:ch21-stab:sum}  Write \(u:=t_1\varphi_1+t_2\varphi_2\), with the
convention that \(t\cdot(-\infty)=-\infty\) for \(t>0\) and \(=0\) for
\(t=0\); no difference of infinities occurs, so \(u\) is well defined with
values in \([-\infty,+\infty)\).  It is upper semicontinuous, being a
nonnegative combination of upper semicontinuous functions.  It satisfies
\eqref{eq:ch21-retraction-inequality}, because
\(\varphi_j\leq\varphi_j\circ p_{\mathcal X}\) and both coefficients are
nonnegative.  It is not identically \(-\infty\): by
\Cref{prop:ch21-first-consequences}\,\cref{itm:ch21-fc:finite} both
\(\varphi_1\) and \(\varphi_2\) are finite on \(\Delta_{\mathcal X}\),
hence so is \(u\).  This last point is not a formality; it is the
reason the theorem is stated for a common determining model, and it is the
first place in the chapter where a sum of two singular potentials is seen
to be nontrivial.

For clause~(2), let \(\psi_{j,k}\) be \(\theta_j\)-psh model functions with
\(\psi_{j,k}|_{\Delta_{\mathcal X}}\to\varphi_j|_{\Delta_{\mathcal X}}\)
uniformly.  Then \(t_1\psi_{1,k}+t_2\psi_{2,k}\) is a model function, and
its curvature form is
\[
 t_1\theta_1+t_2\theta_2
 +\ddc(t_1\psi_{1,k}+t_2\psi_{2,k})
 =t_1(\theta_1+\ddc\psi_{1,k})+t_2(\theta_2+\ddc\psi_{2,k}),
\]
by linearity of \(\ddc\) on \(\mathcal D(X)_{\mathbf R}\)
\eqref{eq:ch20-ddc}, and this is semipositive because the semipositive
forms constitute a convex cone
(\Cref{prop:ch20-semipositive-well-defined}).  So
\(t_1\psi_{1,k}+t_2\psi_{2,k}\) is a
\((t_1\theta_1+t_2\theta_2)\)-psh model function, and
\[
 \bigl\lVert (t_1\psi_{1,k}+t_2\psi_{2,k})-u\bigr\rVert_{\Delta_{\mathcal X}}
 \leq t_1\lVert\psi_{1,k}-\varphi_1\rVert_{\Delta_{\mathcal X}}
     +t_2\lVert\psi_{2,k}-\varphi_2\rVert_{\Delta_{\mathcal X}}
 \longrightarrow0 .
\]

\cref{itm:ch21-stab:max}  Upper semicontinuity of a maximum of two upper
semicontinuous functions is
\Cref{lem:usc-star-calculus}\,(4).  The retraction inequality follows from
\[
 \max\{\varphi,\psi\}
 \leq\max\{\varphi\circ p_{\mathcal X},\psi\circ p_{\mathcal X}\}
 =\max\{\varphi,\psi\}\circ p_{\mathcal X}.
\]
The maximum is not identically \(-\infty\) since \(\varphi\) is not.  For
clause~(2), if \(\psi_k,\chi_k\) are \(\theta_1\)-psh model functions
approximating \(\varphi,\psi\) uniformly on \(\Delta_{\mathcal X}\), then
\(\max\{\psi_k,\chi_k\}\) is again a \(\theta_1\)-psh model function by
\Cref{prop:ch20-model-max}, and the elementary inequality
\[
 \bigl|\max\{a,b\}-\max\{a',b'\}\bigr|
 \leq\max\{|a-a'|,|b-b'|\}
\]
shows that the approximation is again uniform.  Finitely many arguments
follow by induction.

\cref{itm:ch21-stab:unif}  Upper semicontinuity and
\eqref{eq:ch21-retraction-inequality} pass to uniform limits.  On the fixed
complex \(\Delta_{\mathcal X}\), each restriction
\(\varphi_k|_{\Delta_{\mathcal X}}\) is finite and continuous by
\Cref{prop:ch21-first-consequences}, and the restrictions converge uniformly
to \(\varphi|_{\Delta_{\mathcal X}}\).  Thus the latter is finite and
continuous, in particular \(\varphi\not\equiv-\infty\).  For clause~(2), a
diagonal choice suffices: given \(\varepsilon>0\), pick \(k\) with
\(\lVert\varphi_k-\varphi\rVert_{X^{\An}}\leq\varepsilon/2\) and then a
\(\theta_1\)-psh model function \(\psi\) with
\(\lVert\psi-\varphi_k\rVert_{\Delta_{\mathcal X}}\leq\varepsilon/2\).
This proves membership in \(\PSH(X,\theta_1)\); it gives no continuity away
from the dual complexes.

\cref{itm:ch21-stab:mono}  Let \(\psi_k\) be \(\theta_1\)-psh model functions
approximating \(\varphi\) uniformly on \(\Delta_{\mathcal X}\).  Then
\(\theta_2+\ddc\psi_k=(\theta_2-\theta_1)+(\theta_1+\ddc\psi_k)\) is a sum
of two semipositive forms, hence semipositive; so each \(\psi_k\) is also
\(\theta_2\)-psh.  Clause~(1) and upper semicontinuity do not refer to
\(\theta\).
\end{proof}

\begin{remark}[What is \emph{not} in the list]
\label{rem:ch21-not-stable}
Three operations are conspicuously absent, and their absence is deliberate.

\emph{Minima.}  The minimum of two \(\theta\)-psh functions is essentially
never \(\theta\)-psh.  Already for model functions on a curve,
\(\min\{\varphi_1,\varphi_2\}\) creates a concave kink, and
\eqref{eq:ch21-curve-psh-test} fails there;
\Cref{ex:ch21-conic-numbers} computes the resulting negative mass
explicitly.  This is the same asymmetry between \(\sup\) and \(\inf\)
recorded in \Cref{lem:usc-inf} and in the discussion following it.

\emph{Arbitrary suprema.}  A supremum of a family of \(\theta\)-psh
functions, uniformly bounded above, has an upper semicontinuous
regularization that is again \(\theta\)-psh; but the raw supremum need not
be upper semicontinuous, and the proof of the regularized statement needs
the equicontinuity estimate of \Cref{chap:bfj-compactness}.  The statement
is therefore proved as \Cref{thm:ch23-upper-envelope}, not here.  A reader
who wants it now should note that the finite-maximum case
\cref{itm:ch21-stab:max} above is genuinely elementary, whereas the infinite
case is not.

\emph{Products and general compositions.}  There is no analogue of
\Cref{prop:ch15-composition}: composing \(\varphi\) with a convex increasing
function does not in general preserve \(\PSH(X,\theta)\), because the
curvature form of the composite is not a nonnegative combination of
available forms.  The only composition available at model level is the
maximum, and it is available because of the ideal-theoretic mechanism of
\Cref{prop:ch20-model-max}.
\end{remark}

\subsection{Decreasing nets}
\label{sec:ch21-decreasing}

\begin{theorem}\label{thm:ch21-decreasing-stability}
Let \(\varphi_i\) be a decreasing net in \(\PSH(X,\theta)\), and put
\(\varphi=\inf_i\varphi_i\).  Then either
\(\varphi\equiv-\infty\) or \(\varphi\in\PSH(X,\theta)\).
\end{theorem}

\begin{proof}
\emph{Step 1: the two formal clauses.}
A decreasing infimum of upper semicontinuous functions is upper
semicontinuous by \Cref{lem:usc-inf}.  The retraction inequality
\eqref{eq:ch21-retraction-inequality} passes to the
limit: for each \(i\) and each \(x\), \(\varphi(x)\leq\varphi_i(x)\leq
\varphi_i(p_{\mathcal X}(x))\), and taking the infimum over \(i\) on the
right gives \(\varphi(x)\leq\varphi(p_{\mathcal X}(x))\).

\emph{Step 2: two-sided bounds on one dual complex.}
Fix an SNC model \(\mathcal X\) determining \(\theta\).  The restrictions
\(\varphi_i|_{\Delta_{\mathcal X}}\) are continuous and belong to the
common equicontinuous family of \Cref{thm:bfj-skeletal-equicontinuity}: by
part~(1) of that theorem they share a Lipschitz constant \(C\) on each face
of \(\Delta_{\mathcal X}\), where \(C\) depends only on
\((\mathcal X,\theta)\).  Suppose
the global limit is not identically \(-\infty\), and choose
\(x_0\in X^{\An}\) with \(\varphi(x_0)>-\infty\).  Put
\(y_0=p_{\mathcal X}(x_0)\).  The retraction inequality gives

\[
        \varphi_i(y_0)\geq\varphi_i(x_0)\geq\varphi(x_0)
        \quad\text{for every }i,
\]

a lower bound at one point of \(\Delta_{\mathcal X}\) that is uniform in
\(i\).  Fixing one index \(i_0\) and discarding the indices not beyond it,
monotonicity also bounds
\(\varphi_i(y_0)\leq\varphi_{i_0}(y_0)\) from above by one fixed value.
The common Lipschitz constant, together with the finite diameter of the
compact polyhedron \(\Delta_{\mathcal X}\), therefore gives two-sided
uniform bounds
\[
 \varphi(x_0)-C\operatorname{diam}\Delta_{\mathcal X}
 \ \leq\ \varphi_i\ \leq\
 \varphi_{i_0}(y_0)+C\operatorname{diam}\Delta_{\mathcal X}
 \qquad\text{on }\Delta_{\mathcal X},
\]
valid for all \(i\geq i_0\).

\emph{Step 3: uniform convergence on the dual complex.}
The pointwise limit \(\varphi|_{\Delta_{\mathcal X}}\) is therefore finite,
and it inherits the Lipschitz constant \(C\) on each face, hence is
continuous.  The net version of Dini's theorem
\Cref{prop:dini-nets}, applied to the decreasing net of continuous
nonnegative functions
\(\varphi_i|_{\Delta_{\mathcal X}}-
\varphi|_{\Delta_{\mathcal X}}\) on the compact space
\(\Delta_{\mathcal X}\), now gives uniform convergence.  Equivalently one
may invoke \Cref{prop:dini-usc}\,(1) directly, whose hypothesis is exactly
that the limit is real valued and lower semicontinuous.

\emph{Step 4: clause \textup{(2)}.}
Each \(\varphi_i|_{\Delta_{\mathcal X}}\) lies in the uniformly closed set
\(\mathcal P_{\mathcal X}(\theta)\) of \eqref{eq:ch21-skeletal-closure}, and
by Step~3 the limit is a uniform limit of these; a uniformly closed set
contains its uniform limits, so
\(\varphi|_{\Delta_{\mathcal X}}\in\mathcal P_{\mathcal X}(\theta)\), which
is clause~(2).  Since \(\mathcal X\) was an arbitrary determining model,
\(\varphi\in\PSH(X,\theta)\).

This is \cite[Corollary~7.9]{BFJ16}; the uniform
skeletal estimate, proved in \Cref{chap:bfj-compactness}, is the
nonformal input.
\end{proof}

\begin{remark}[Where the net is essential, and where it is not]
\label{rem:ch21-decreasing-nets}
The use of a net is essential in the statement: the directed system of
models has no countable cofinal subset in general
(\Cref{lem:ch21-no-cofinal-sequence}), the Berkovich space \(X^{\An}\) is
not metrizable in general, and \Cref{thm:bfj-decreasing-regularization}
produces exactly a net.  \Cref{ex:decreasing-net-not-sequential} exhibits a
decreasing net of upper semicontinuous functions on a compact Hausdorff
space that admits no decreasing subsequence with the same infimum, which is
the abstract shape of the difficulty.

Inside the proof, by contrast, the net plays almost no role.  Steps~1
and~4 are purely order-theoretic.  Step~2 uses only that the index set is
directed, so as to discard an initial segment.  Step~3 uses Dini's theorem
in the net form \Cref{prop:dini-nets}, which is proved for nets precisely
because directedness, and not countability, is what the finite-subcover
argument needs.  This division is typical: the statements of
\Cref{part:bfj} require nets, and the proofs require nothing more than
directedness.
\end{remark}

\subsection{Pullback}
\label{sec:ch21-pullback}

Pullback is the one operation in the plan for this chapter whose general
form is genuinely deferred.  We prove it here on the two loci where it can
be proved with the tools of \Cref{chap:bfj-model-forms}, namely for model
functions and for continuous psh functions, and we say precisely what is
missing in the singular case.

Let \(f:Y\to X\) be a morphism of smooth connected projective
\(K\)-varieties.  If \(\theta\in\mathcal Z^{1,1}(X)\) is determined on
\(\mathcal X\), choose a model \(\mathcal Y\) of \(Y\) together with a
morphism \(\tilde f:\mathcal Y\to\mathcal X\) extending \(f\): take the
closure of the graph of \(f\) inside \(\mathcal Y_0\times_R\mathcal X\) for
any model \(\mathcal Y_0\) of \(Y\), normalize, and resolve.  Set

\begin{equation}
\label{eq:ch21-pullback-form}
 f^*\theta:=\bigl[\tilde f^*\theta_{\mathcal X}\bigr]
 \in\mathcal Z^{1,1}(Y).
\end{equation}

This is independent of all choices, because two such extensions agree after
passing to a model of \(Y\) dominating both.

\begin{proposition}
\label{prop:ch21-pullback}
With the notation above:
\begin{enumerate}
\item\label{itm:ch21-pb:model}
if \(\varphi\) is a \(\theta\)-psh model function on \(X\), then
\(f^*\varphi:=\varphi\circ f^{\An}\) is an \(f^*\theta\)-psh model function
on \(Y\);
\item\label{itm:ch21-pb:cont}
if \(\varphi\in\PSH(X,\theta)\) is continuous, then
\(f^*\varphi\in\PSH(Y,f^*\theta)\) and is continuous;
\item\label{itm:ch21-pb:sup}
in both cases \(\sup_{Y^{\An}}f^*\varphi\leq\sup_{X^{\An}}\varphi\), with
equality if \(f\) is surjective.
\end{enumerate}
\end{proposition}

\begin{proof}
\cref{itm:ch21-pb:model}  Choose a model \(\mathcal X'\geq\mathcal X\) on
which \(\varphi=\varphi_D\) for a vertical \(\mathbf R\)-Cartier divisor
\(D\), and a morphism \(\tilde f:\mathcal Y\to\mathcal X'\) as above.  The
pullback \(\tilde f^*D\) is defined because \(D\) is Cartier, and it is
vertical because \(\tilde f^{-1}(\mathcal X'_s)=\mathcal Y_s\), both being
the fibres over the closed point of \(\Spec R\).  Its local equations pull
back, so \(\varphi_{\tilde f^*D}=\varphi_D\circ f^{\An}\) by the sign
convention \eqref{eq:ch20-sign}; in particular \(f^*\varphi\) is a model
function.  Its curvature form is
\[
 f^*\theta+\ddc f^*\varphi
 =\bigl[\tilde f^*\theta_{\mathcal X'}+\tilde f^*D\bigr]
 =\bigl[\tilde f^*(\theta_{\mathcal X'}+D)\bigr],
\]
which is semipositive because the pullback of a relatively nef class along a
morphism of models is relatively nef
(\Cref{prop:ch20-semipositive-well-defined}, \Cref{prop:nef-pullback}).

\cref{itm:ch21-pb:cont}  By \Cref{thm:ch21-continuous-zhang} a continuous
\(\theta\)-psh function is a uniform limit on \(X^{\An}\) of a sequence of
\(\theta\)-psh model functions \(\psi_k\).  Since \(f^{\An}\) is continuous,
\(f^*\psi_k\to f^*\varphi\) uniformly on \(Y^{\An}\), and each
\(f^*\psi_k\) is an \(f^*\theta\)-psh model function by
\cref{itm:ch21-pb:model}.  Now apply
\Cref{thm:ch21-stability}\,\cref{itm:ch21-stab:unif} on \(Y\).

\cref{itm:ch21-pb:sup}  Immediate, since \(f^*\varphi\) takes values in the
set of values of \(\varphi\) on \(f^{\An}(Y^{\An})\), and \(f^{\An}\) is
surjective when \(f\) is, by properness.
\end{proof}

\begin{warning}{Pullback of a singular potential is a later theorem}
\Cref{prop:ch21-pullback} is not asserted for a general
\(\varphi\in\PSH(X,\theta)\), and the reader should not assume it.  The
obstruction is visible in the proof: the argument passes through uniform
approximation on all of \(X^{\An}\), which is available only in the
continuous case, and the alternative route through
\eqref{eq:ch21-skeletal-reconstruction} fails because retractions do not
commute with \(f\): in general
\(p_{\mathcal X}\circ f^{\An}\neq f^{\An}\circ p_{\mathcal Y}\), so
\(f^*(\varphi\circ p_{\mathcal X})\) is not \((f^*\varphi)\circ
p_{\mathcal Y}\).  There is moreover a genuine issue of definition when
\(f(Y)\) is contained in the polar set of \(\varphi\), in which case
\(f^*\varphi\equiv-\infty\) and \Cref{def:ch21-bfj-psh} excludes it; this is
the BFJ counterpart of item (L8) of \Cref{sec:ch19-limitations-table}.  The
general statement, with the hypotheses that make it true --- surjectivity or
a condition on the image relative to the polar set, and a separate treatment
of ground-field extension --- is
\Cref{chap:bfj-functoriality}.
\end{warning}

\begin{remark}[Ground-field extension]
\label{rem:ch21-base-change}
The pullback along a complete extension \(K'/K\) is a different operation
from \Cref{prop:ch21-pullback} and is also deferred.  It is well behaved for
model functions, since a model base-changes and relative nefness may be
tested after a faithfully flat base change of the valuation ring, which is
the mechanism used in Step~1 of the proof of
\Cref{thm:ch19-cld-bfj-comparison}.  What is not formal is that the
extension \(K'\) need not be discretely valued --- already
\(K'=\widehat{K^a}\) is not --- so the target of the base change lies
outside \eqref{eq:ch21-hypotheses} and outside the framework of the present
chapter.  \Cref{chap:bfj-functoriality} treats this and records which
statements survive.
\end{remark}

\section{Singularities and polar sets}
\label{sec:ch21-singularities}

\subsection{Singularity types}
\label{sec:ch21-singularity-types}

\begin{definition}
\label{def:ch21-polar-singularity}
\index{polar set!BFJ}\index{singularity type!BFJ}
The \emph{polar set} of \(\varphi\in\PSH(X,\theta)\) is

\[
                         \{\varphi=-\infty\}.
\]

For \(\varphi,\psi\in\PSH(X,\theta)\) we write \(\varphi\preceq\psi\), and
say that \(\varphi\) is \emph{more singular} than \(\psi\), if
\(\varphi\leq\psi+C\) for some \(C\in\mathbf R\).  Two psh functions have
the same \emph{singularity type} if
\(\varphi\preceq\psi\preceq\varphi\), equivalently if their difference
is bounded in both directions, equivalently if there is \(C\geq0\) such that

\begin{equation}
\label{eq:ch21-singularity-type}
                         \varphi-C\leq\psi\leq\varphi+C.
\end{equation}
\end{definition}

The relation \(\preceq\) is reflexive and transitive, so singularity types
form a partially ordered set.  Note that the definition compares two
elements of the \emph{same} class \(\PSH(X,\theta)\); the difference of two
such functions is not \(\theta\)-psh for any form and is not even defined
where both are \(-\infty\), so \eqref{eq:ch21-singularity-type} is to be
read as the pair of inequalities displayed, each of which is meaningful in
\([-\infty,+\infty)\) by the conventions of
\Cref{sec:extended-arithmetic}.

\begin{proposition}
\label{prop:ch21-singularity-elementary}
Let \(\varphi,\psi\in\PSH(X,\theta)\).
\begin{enumerate}
\item\label{itm:ch21-sing:polar}
If \(\varphi\preceq\psi\), then
\(\{\psi=-\infty\}\subseteq\{\varphi=-\infty\}\).  In particular functions
of the same singularity type have the same polar set.
\item\label{itm:ch21-sing:bounded}
All bounded members of \(\PSH(X,\theta)\) have the same singularity type,
namely that of the constants when \(\theta\) is semipositive, and their
common polar set is empty.
\item\label{itm:ch21-sing:max}
\(\max\{\varphi,\psi\}\) is the least upper bound of \(\varphi\) and
\(\psi\) for \(\preceq\) inside \(\PSH(X,\theta)\), and its polar set is
\(\{\varphi=-\infty\}\cap\{\psi=-\infty\}\).
\item\label{itm:ch21-sing:convex}
For \(t\in(0,1)\) the polar set of \(t\varphi+(1-t)\psi\) is
\(\{\varphi=-\infty\}\cup\{\psi=-\infty\}\), and its singularity type
depends only on those of \(\varphi\) and \(\psi\).
\item\label{itm:ch21-sing:divisorial}
The relation \(\varphi\preceq\psi\) may be tested at divisorial points
alone.
\end{enumerate}
\end{proposition}

\begin{proof}
\cref{itm:ch21-sing:polar}  If \(\psi(x)=-\infty\) and
\(\varphi\leq\psi+C\), then \(\varphi(x)=-\infty\).

\cref{itm:ch21-sing:bounded}  If \(\varphi,\psi\) are both bounded, then
\(\varphi-\psi\) is a bounded real-valued function and
\eqref{eq:ch21-singularity-type} holds with \(C=\lVert\varphi\rVert_\infty
+\lVert\psi\rVert_\infty\).  Every member of \(\PSH(X,\theta)\) is bounded
above by \Cref{prop:ch21-sup-on-skeleton}, so ``bounded'' means ``bounded
below''.  When \(\theta\) is semipositive the constants belong to the class
by \Cref{prop:ch21-constants}\,\cref{itm:ch21-const:const}.

\cref{itm:ch21-sing:max}  It is an upper bound because
\(\varphi,\psi\leq\max\{\varphi,\psi\}\), and it lies in the class by
\Cref{thm:ch21-stability}\,\cref{itm:ch21-stab:max}.  If \(\chi\) is another
upper bound, say \(\varphi\leq\chi+C_1\) and \(\psi\leq\chi+C_2\), then
\(\max\{\varphi,\psi\}\leq\chi+\max\{C_1,C_2\}\).  The polar statement is
the identity \(\max\{a,b\}=-\infty\) if and only if \(a=b=-\infty\).

\cref{itm:ch21-sing:convex}  For \(t\in(0,1)\), \(t a+(1-t)b=-\infty\) if and
only if \(a=-\infty\) or \(b=-\infty\).  If \(\varphi'\) has the same type
as \(\varphi\) and \(\psi'\) the same as \(\psi\), the constants add:
\(|t\varphi+(1-t)\psi-t\varphi'-(1-t)\psi'|\leq tC_\varphi+(1-t)C_\psi\)
wherever all four are finite, and elsewhere both sides are \(-\infty\)
simultaneously by the previous sentence.

\cref{itm:ch21-sing:divisorial}  Apply
\Cref{prop:ch21-divisorial-order} to \(\varphi\) and \(\psi+C\), the latter
lying in \(\PSH(X,\theta)\) by
\Cref{prop:ch21-constants}\,\cref{itm:ch21-const:transl}.
\end{proof}

\subsection{Polar sets}
\label{sec:ch21-polar}

Every psh function is finite at every quasimonomial point, since such a
point lies on a dual complex.  Thus polar sets are disjoint from the union
of all skeleta, although that union is dense.  This is possible because a
dense subset need not be closed and a singular psh function need not be
continuous globally.  The following proposition collects what is true about
polar sets at this stage of the theory, and the remark after it records what
is not.

\begin{proposition}
\label{prop:ch21-polar-structure}
Let \(\varphi\in\PSH(X,\theta)\) and put \(P(\varphi)=\{\varphi=-\infty\}\).
\begin{enumerate}
\item\label{itm:ch21-pol:gdelta}
\(P(\varphi)=\bigcap_{k\geq1}\{\varphi<-k\}\) is a \(G_\delta\) subset of
\(X^{\An}\).
\item\label{itm:ch21-pol:skeleta}
\(P(\varphi)\) is disjoint from \(\bigcup_{\mathcal X}\Delta_{\mathcal X}\);
equivalently it contains no quasimonomial point, hence no divisorial point.
\item\label{itm:ch21-pol:interior}
\(P(\varphi)\) has empty interior, and its complement is dense.
\item\label{itm:ch21-pol:notclosed}
\(P(\varphi)\) need not be closed, and it need not be Zariski-closed; see
\Cref{ex:ch21-polar-noncompact}.
\item\label{itm:ch21-pol:sigma}
\(P(\max\{\varphi,\psi\})=P(\varphi)\cap P(\psi)\) and
\(P(t\varphi+(1-t)\psi)=P(\varphi)\cup P(\psi)\) for \(0<t<1\); in
particular the polar sets of the class are stable under finite unions and
finite intersections.
\item\label{itm:ch21-pol:curve}
If \(\dim X=1\), then \(P(\varphi)\) consists of points of type \(1\) and
type \(4\) only.
\end{enumerate}
\end{proposition}

\begin{proof}
\cref{itm:ch21-pol:gdelta}  Each \(\{\varphi<-k\}\) is open by upper
semicontinuity, and the intersection over \(k\) is the polar set.

\cref{itm:ch21-pol:skeleta} and \cref{itm:ch21-pol:interior} are
\Cref{prop:ch21-first-consequences},
parts \cref{itm:ch21-fc:quasimonomial} and
\cref{itm:ch21-fc:emptyinterior}.

\cref{itm:ch21-pol:notclosed}  See the example cited.

\cref{itm:ch21-pol:sigma}  This is
\Cref{prop:ch21-singularity-elementary},
parts \cref{itm:ch21-sing:max} and \cref{itm:ch21-sing:convex}; both
operations stay inside \(\PSH(X,\theta)\) by
\Cref{thm:ch21-stability}.

\cref{itm:ch21-pol:curve}  On a smooth projective curve the quasimonomial
points are exactly the points of type \(2\) and type \(3\), and every such
point lies on the skeleton of some SNC model---semistable models would not
suffice, by \Cref{rem:ch21-semistable-not-cofinal}---(\Cref{rem:quasimonomial-vs-divisorial},
\Cref{thm:berk-a1-classification} for \(\mathbf P^1\) and the
corresponding classification for a general curve through
\Cref{prop:ch08-annulus-skeleton}).  Now apply
\cref{itm:ch21-pol:skeleta}.
\end{proof}

\begin{warning}{Two things a polar set is not}
First, a polar set here is \emph{not} the same notion as the pluripolar sets
of \Cref{chap:bfj-capacity}, which are defined by a capacity and form a
\(\sigma\)-ideal.  \Cref{prop:ch21-polar-structure}\,\cref{itm:ch21-pol:sigma}
gives only finite stability; countable stability requires the capacity
estimates of \Cref{prop:polar-sigma-ideal} and is proved there, not here.

Second, no analyticity is asserted.  \Cref{ex:ch21-polar-noncompact}
produces a psh function on \(\mathbf P^1\) whose polar set is a countably
infinite set of rational points, which is not Zariski-closed; this is the
exact analogue of \Cref{ex:ch15-polar-nonanalytic} in the CLD setting, where
the same warning is issued.  A Siu-type analyticity theorem for polar sets
or for Lelong sublevel sets is not available in this theory, in either the
CLD or the BFJ formulation.
\end{warning}

\begin{remark}[Polar sets and the Monge--Amp\`ere operator]
\label{rem:ch21-polar-ma}
It is tempting to say that a psh function with nonempty polar set has a
curvature ``measure'' charging that set, on the model of
\(\theta+\ddc\log\lVert s\rVert_0=\delta_{\divis(s)}\) for a section \(s\).
No such statement is made in this chapter, and the reader should note why.
For a general \(\varphi\in\PSH(X,\theta)\) no measure has yet been attached
to \(\theta+\ddc\varphi\): the monotone extension of the mixed model measure
to bounded psh functions is \Cref{thm:bounded-ma-extension}, and the
non-pluripolar product is \Cref{def:nonpluripolar-ma}, both in later
chapters, and the two disagree on polar sets by construction.  What
\emph{is} legitimate here, and what the examples of
\Cref{sec:ch21-worked} do, is to compute the model measures of the
truncations \(\max\{\varphi,-m\}\), each of which is a model function, and
to observe where their mass goes as \(m\to\infty\).
\end{remark}

\section{Comparison with continuous semipositive metrics}
\label{sec:ch21-continuous}

For continuous metrics the BFJ and Zhang definitions agree.

\begin{theorem}\label{thm:ch21-continuous-zhang}
Let \(L\) be ample and fix a model reference metric of curvature \(\theta\).
For a continuous function \(\varphi\), the following are equivalent:

\begin{enumerate}
\item \(\varphi\in\PSH(X,\theta)\);
\item \(\varphi\) is a uniform limit on \(X^{\An}\) of \(\theta\)-psh model
functions;
\item the metric \(\|\cdot\|_0\rme^{-\varphi}\) is semipositive in Zhang's
sense.
\end{enumerate}
\end{theorem}

\begin{proof}
The equivalence of (2) and (3) is Zhang's definition, recalled in
\Cref{sec:metrics-semipositive}: a continuous metric on an ample \(L\) is
semipositive in Zhang's sense precisely when it is a uniform limit of
semipositive model metrics, and by
\Cref{cor:ch21-reference-independence} the relative weights of those model
metrics are exactly the \(\theta\)-psh model functions.

(2)\(\Rightarrow\)(1).  Uniform convergence on \(X^{\An}\)
immediately implies the two clauses of \Cref{def:ch21-bfj-psh}: clause~(2)
because restriction to \(\Delta_{\mathcal X}\) does not increase the
sup-norm, and clause~(1) because each model psh function satisfies it by
\Cref{lem:ch21-model-skeletal-properties} and inequalities pass to uniform
limits.  This implication is
\Cref{thm:ch21-stability}\,\cref{itm:ch21-stab:unif}.

(1)\(\Rightarrow\)(2).  The decreasing regularization theorem
\Cref{thm:bfj-decreasing-regularization} supplies a decreasing net
\((\psi_i)\) of \(\theta\)-psh model functions with
\(\psi_i\searrow\varphi\) pointwise on \(X^{\An}\).  Each \(\psi_i\) is
continuous, and by hypothesis the limit \(\varphi\) is continuous, so
Dini's theorem for nets on the compact space
\(X^{\An}\), \Cref{prop:dini-nets}, gives \(\psi_i\to\varphi\) uniformly;
the same conclusion follows from \Cref{prop:dini-usc}\,(1) or from the
appendix form \Cref{thm:appE-dini}.  For each \(m\geq1\) select an index
\(i_m\) with \(\lVert\psi_{i_m}-\varphi\rVert_\infty\leq1/m\).  The
resulting \emph{sequence} \((\psi_{i_m})_{m\geq1}\) of \(\theta\)-psh model
functions converges uniformly to \(\varphi\), which is (2).  The algebraic
proof of decreasing regularization is completed in
\Cref{chap:bfj-envelopes}.
\end{proof}

\begin{remark}[A legitimate passage from a net to a sequence]
\label{rem:ch21-net-to-sequence}
The last proof contains the one place in this chapter where a net is
replaced by a sequence, and it is instructive to see why this is allowed.
The replacement is not a cofinality statement about the index set --- no
such statement is available, by \Cref{lem:ch21-no-cofinal-sequence}.  It is
a statement about a \emph{metric}: continuity of the limit converts the
decreasing net into a net convergent in the Banach space
\(C^0(X^{\An})\), and in a metric space one may always extract a sequence
realizing a given sequence of accuracies.  The extracted sequence is not
increasing or decreasing in the index, and it is not cofinal; it is only
uniformly convergent, which is all that (2) asks for.  Compare the
checklist \Cref{rem:net-vs-sequence-checklist}: this is an instance of the
first mechanism, metrizability, applied to the target rather than to the
source.
\end{remark}

\begin{corollary}
\label{cor:ch21-continuous-cone}
The set \(\PSH(X,\theta)\cap C^0(X^{\An})\) is convex, stable under finite
maxima and under addition of constants, and closed in \(C^0(X^{\An})\) for
the uniform norm.  It is the uniform closure of the set of \(\theta\)-psh
model functions.
\end{corollary}

\begin{proof}
Convexity and stability are
\Cref{thm:ch21-stability}, parts \cref{itm:ch21-stab:sum},
\cref{itm:ch21-stab:max} and \cref{itm:ch21-stab:const}, since all the
operations preserve continuity.  Closedness is
\Cref{thm:ch21-stability}\,\cref{itm:ch21-stab:unif}.  The last assertion is
the equivalence of (1) and (2) in \Cref{thm:ch21-continuous-zhang}.
\end{proof}

\subsection{Exactly what is being compared, and what is not}
\label{sec:ch21-cld-contrast}

\Cref{thm:ch21-continuous-zhang} concerns continuous global metrics.  It
does not identify
arbitrary BFJ singular psh functions with every local CLD or curve-test psh
function; the precise comparison locus is recorded in
\Cref{chap:master-dictionary} and proved in
\Cref{thm:ch19-cld-bfj-comparison}.  This subsection states the scope
exactly, in the form the rest of \Cref{part:bfj} relies on.

\begin{center}
\small
\begin{tabular}{p{.30\textwidth}|p{.30\textwidth}|p{.32\textwidth}}
\toprule
\textbf{Locus}&\textbf{What is proved}&\textbf{Reference}\\
\midrule
model functions
 &\(\theta\)-psh \(\Leftrightarrow\) Zhang semipositive
  \(\Leftrightarrow\) positive Chern current
  \(\Leftrightarrow\) positive Chern delta-form
  \(\Leftrightarrow\) functorially psh weights
  \(\Leftrightarrow\) curve test
 &\Cref{thm:ch19-cld-bfj-comparison}\,(A);
  here \Cref{thm:ch21-model-consistency} for the first equivalence\\
continuous \(\theta\)-psh functions
 &uniform limits of \(\theta\)-psh model functions; the two
  Monge--Amp\`ere measures agree
 &\Cref{thm:ch21-continuous-zhang} and
  \Cref{thm:ch19-cld-bfj-comparison}\,(B)\\
bounded \(\theta\)-psh functions
 &\emph{nothing}: not asserted to be locally CLD-psh
 &\Cref{thm:ch19-cld-bfj-comparison}\,\cref{itm:ch19-cb:scope1}\\
locally CLD-psh functions on \(X^{\An}\)
 &\emph{nothing}: not asserted to be \(\theta\)-psh for any \(\theta\)
 &\Cref{thm:ch19-cld-bfj-comparison}\,\cref{itm:ch19-cb:scope2}\\
unbounded \(\theta\)-psh functions
 &\emph{nothing}: the two Monge--Amp\`ere constructions are not compared,
  and they disagree on polar sets
 &\Cref{thm:ch19-cld-bfj-comparison}\,\cref{itm:ch19-cb:scope3},
  \cref{itm:ch19-cb:scope4}\\
\bottomrule
\end{tabular}
\end{center}

Two structural separations \emph{are} provable, and they are worth having
because they show that the empty entries above are not merely gaps in the
present state of knowledge but reflect a difference of type between the two
constructions.

\begin{proposition}
\label{prop:ch21-bfj-vs-cld}
Assume \eqref{eq:ch21-hypotheses}.
\begin{enumerate}
\item\label{itm:ch21-sep:global}
\emph{The BFJ class is not a sheaf.}  There is no notion of
``\(\theta\)-psh on \(U\)'' for a proper open subset
\(U\subsetneq X^{\An}\): \Cref{def:ch21-bfj-psh} quantifies over dual
complexes of models of \(X\), which are not contained in \(U\).  By
contrast the CLD class of \Cref{def:ch15-hierarchy} is a sheaf, and its
sections over open subsets include functions unbounded above.  Explicitly,
for \(X=\mathbf P^1_K\) and \(U=\mathbf P^{1,\An}_K\setminus\{0\}\), the
function \(u=\log|1/z|\) is a locally approachable CLD-psh function on
\(U\) with \(\sup_Uu=+\infty\), whereas every member of
\(\PSH(\mathbf P^1,\theta)\) satisfies
\(\sup\varphi=\max_{\Delta_{\mathcal X}}\varphi<+\infty\).
\item\label{itm:ch21-sep:nets}
\emph{The two classes are generated by different limit operations.}
\(\PSH(X,\theta)\) is stable under decreasing \emph{nets}
(\Cref{thm:ch21-decreasing-stability}) and every member is the limit of a
decreasing net of model functions
(\Cref{thm:bfj-decreasing-regularization}); the index set of that net has no
countable cofinal subfamily in general
(\Cref{lem:ch21-no-cofinal-sequence}).  The CLD class
\(\operatorname{Psh}_{S,1}\) is generated by decreasing \emph{sequences} and
is not known to be stable under further decreasing limits
(\Cref{thm:ch15-approachable-stability}).
\end{enumerate}
\end{proposition}

\begin{proof}
\cref{itm:ch21-sep:global}  The first sentence is a reading of
\Cref{def:ch21-bfj-psh}: both clauses refer to \(\Delta_{\mathcal X}\) and
to \(p_{\mathcal X}\), which are global objects.  For the example, note that
\(1/z\) is a regular analytic function on
\(U=\mathbf P^{1,\An}_K\setminus\{0\}\), with a single zero at \(\infty\);
by \Cref{prop:ch15-log-approachable} the function \(\log|1/z|\) is a
locally approachable psh function on \(U\), with polar locus
\(\{\infty\}\).  Its supremum is \(+\infty\): along the points
\(\zeta_{0,r}\) with \(r\downarrow0\) one has
\(\log|1/z|(\zeta_{0,r})=-\log r\to+\infty\).  The boundedness of BFJ psh
functions is \Cref{prop:ch21-sup-on-skeleton}.

\cref{itm:ch21-sep:nets}  Each assertion is the cited statement.
\end{proof}

\begin{remark}[Why no separating example on \(X^{\An}\) is available here]
\label{rem:ch21-no-separating-example}
\Cref{prop:ch21-bfj-vs-cld} separates the two classes as \emph{classes of
functions on varying domains}.  It does not produce a function on the whole
of \(X^{\An}\) that belongs to one class and not the other, and this book
contains no such function.  The reason is precise and worth stating, since
it is the exact content of
\Cref{thm:ch19-cld-bfj-comparison}\,\cref{itm:ch19-cb:scope1} and
\cref{itm:ch19-cb:scope2}.

In the direction ``BFJ \(\Rightarrow\) CLD'': a \(\theta\)-psh function is a
decreasing net limit of model functions; each model function has locally
CLD-psh weights, by
\Cref{thm:ch19-cld-bfj-comparison}\,\cref{itm:ch19-cb:funct}; so a proof
would need the CLD class to be stable under decreasing nets of its members.
It is not known to be stable even under decreasing \emph{sequences} beyond
the first stage, which is the open problem recorded at
\Cref{thm:ch15-approachable-stability} and in item (L5) of
\Cref{sec:ch19-limitations-table}.  Producing a counterexample instead would
require exhibiting a bounded \(\theta\)-psh function that fails a local
regularity condition; no method for verifying such a failure is available,
because the only known certificates for membership in the CLD sheaf are
constructive.

In the direction ``CLD \(\Rightarrow\) BFJ'': a locally CLD-psh function on
\(X^{\An}\) is automatically upper semicontinuous and bounded above, since
\(X^{\An}\) is compact, so the cheap obstruction of
\cref{itm:ch21-sep:global} disappears.  What is missing is any mechanism
converting local positivity of a current into the \emph{global} numerical
positivity on models that clause~(2) requires.  The local
current-positive-to-first-stage implication in item (L4) of
\Cref{sec:ch19-limitations-table} is a theorem, in its separate local scope,
for continuous functions on $C^{\An}$ when $K$ is algebraically closed,
complete, and nontrivially valued and $C$ is a smooth proper algebraic curve
(\Cref{thm:p063-current-psh-ordinary-regularization}).  Neither it nor
Wanner's comparison supplies this global numerical conversion, and no
descent of the ordinary-smooth approximants to the standing BFJ base field
is asserted.  The book proves the comparison on the model and continuous
metric loci in all dimensions and asserts nothing further.
\end{remark}

\section{Worked examples}
\label{sec:ch21-worked}

Seven computations follow, in six groups.  The first is the elementary kink and fixes the
signs.  The second is a psh function on a curve skeleton whose restrictions
to successive dual complexes are computed as the model is refined, with all
masses checked against the degree.  The third computes a maximum and a
convex combination on that same skeleton, to explicit numbers.  The fourth
computes a polar set, and exhibits a psh function whose polar set is
countably infinite and not closed.  The fifth is the toric picture, and the
sixth is the change of reference form.

All computations use the curve conventions of \Cref{chap:model-calculations}:
the graph Laplacian \eqref{eq:graph-laplacian} is the sum of outgoing
slopes, so a convex kink carries positive mass; the curvature measure is
\eqref{eq:graph-curvature}; and the total mass of
\(\mu_0+\Delta_\Gamma\varphi\) equals \(\deg_XL\) by
\eqref{eq:lap-mass-zero}.  Every displayed measure below is checked against
that degree.  A reader whose masses disagree has almost certainly used the
opposite sign in \eqref{eq:ch20-sign}.

\subsection{A kink, and the signs}
\label{sec:ch21-worked-kink}

\begin{example}[A kink on a semistable skeleton]
\label{ex:ch21-kink-example}
On the nodal interval of \Cref{chap:skeleta}, let

\begin{equation}
\label{eq:ch21-kink}
                  u(t)=\max\{-pt,-q(1-t)\},
                  \qquad 0\leq t\leq1.
\end{equation}

By \eqref{eq:ch20-nodal-ideal}, this is the model function of
\((x^p,y^q)\).  It is convex, its left slope is \(-p\), its right slope is
\(q\), and its slope jumps upward by \(p+q\) at \(q/(p+q)\).  If a reference
form supplies enough endpoint degree to make the two endpoint inequalities
\eqref{eq:ch20-curve-nef-test} nonnegative, then \(u\) is
\(\theta\)-psh.  Refining the model merely subdivides the same convex graph.

Explicitly, with \(\Gamma=[0,1]\) and a reference curvature measure
\(\mu_0=d_0\delta_0+d_1\delta_1\) of total mass \(d=d_0+d_1\),
\eqref{eq:graph-laplacian} gives
\[
 \Delta_\Gamma u
 =-p\,\delta_0+(p+q)\,\delta_{q/(p+q)}-q\,\delta_1 ,
\]
so that
\[
 \mu_0+\Delta_\Gamma u
 =(d_0-p)\delta_0+(p+q)\delta_{q/(p+q)}+(d_1-q)\delta_1 ,
\]
of total mass \(d_0+d_1=d\), as \eqref{eq:lap-mass-zero} requires.  Thus
\(u\) is the restriction of a \(\theta\)-psh model function exactly when
\(p\leq d_0\) and \(q\leq d_1\).  Three readings are worth recording.  The
slope bounds are \emph{two-sided}: convexity alone does not suffice, and it
is the endpoint inequalities, that is the numerical positivity at the two
components of the special fibre, that limit how steep a psh function may be.
Refining the model by a toroidal blowup subdivides \([0,1]\) at rational
points and changes nothing, in accordance with
\Cref{rem:ch21-cofinal-families}\,(c).  And the mass \(p+q\) at the kink is
exactly the jump of the slope, which is the one-dimensional shadow of the
general principle that curvature measures the failure of a potential to be
affine.
\end{example}

\subsection{A psh function on a curve skeleton, refined}
\label{sec:ch21-worked-conic}

The next example is the promised computation of the restrictions of one psh
function to a tower of dual complexes.  We use the running conic of
\Cref{sec:running-semistable}, so that all the intersection numbers have
already been computed in \Cref{chap:model-calculations}.

\begin{example}[A section of \(\cO(1)\) on the conic, and its skeletal
shadows]
\label{ex:ch21-conic-refinement}
\index{plurisubharmonic function!BFJ}\index{polar set!BFJ}
Let \(K\) be complete discretely valued with \(v(\pi)=1\), let
\[
 \mathcal C=\{XY=\pi Z^2\}\subset\mathbf P^2_R
\]
be the model \eqref{eq:running-conic}, with smooth generic fibre the conic
\(C\) and special fibre the union of the two lines
\(C_0=\{Y=0\}\) and \(C_1=\{X=0\}\), meeting at the node
\([0:0:1]\).  Let \(\mathcal L=\cO_{\mathbf P^2_R}(1)|_{\mathcal C}\),
\(L=\mathcal L|_C\) of degree \(2\), and let \(\|\cdot\|_0\) be the
associated model metric, of curvature form \(\theta\).  The skeleton is
\(\Gamma_0=\Delta_{\mathcal C}=[0,1]\), parameterized by
\(s=v(x)\) with \(x=X/Z\), so that \(v(y)=1-s\) for \(y=Y/Z\); the
endpoints are \(v_0=x_{C_0}\) at \(s=0\) and \(v_1=x_{C_1}\) at \(s=1\).
By \Cref{prop:running-conic-measure} the curvature measure is
\begin{equation}
\label{eq:ch21-conic-mu0}
 \mu_0=\delta_{v_0}+\delta_{v_1},
 \qquad
 \text{mass }2=\deg_CL .
\end{equation}

\emph{The potential.}  Take the section \(s_Z:=Z\) of \(L\) and put
\[
 \varphi:=\log\lVert Z\rVert_0
 =\log\frac{|Z|}{\max\{|X|,|Y|,|Z|\}} .
\]
Then \(\varphi\leq0\), and \(\varphi\) is upper semicontinuous with values
in \([-\infty,0]\).  Its divisor is
\(\divis(Z)|_C=\xi_0+\xi_1\) with
\[
 \xi_0=[1:0:0],\qquad \xi_1=[0:1:0],
\]
two \(K\)-rational points of \(C\), because \(Z=0\) and \(XY=\pi Z^2\) force
\(XY=0\).

\emph{Restriction to \(\Gamma_0\).}  Normalizing \(Z=1\), a point
\(\gamma_s\in\Gamma_0\) has \(v(x)=s\), \(v(y)=1-s\), \(v(1)=0\), so
\(\max\{|X|,|Y|,|Z|\}=1\) for every \(s\in[0,1]\) and
\begin{equation}
\label{eq:ch21-conic-restriction-0}
 \varphi|_{\Gamma_0}\equiv0 .
\end{equation}
At the coarsest model the potential is therefore indistinguishable from the
reference, and \(\theta+\ddc(\varphi\circ p_{\mathcal C})=\mu_0\).

\emph{One blowup.}  The point \(\xi_0=[1:0:0]\) reduces to a smooth point
of \(C_0\), namely \(\overline{[1:0:0]}\); indeed in the chart \(X=1\) with
coordinates \(u=Y/X\) and \(w=Z/X\) the equation becomes \(u=\pi w^2\), so
\(\mathcal C\) is locally \(\Spec R[w]\) and \(\xi_0\) is \(w=0\).  Let
\(\rho_1:\mathcal C_1\to\mathcal C\) be the blowup of the closed point
\(\overline{w=0}\) of the special fibre.  The exceptional component \(E\) has
multiplicity one, and the local equation at the new node is again
\(xy=\pi\), so the new edge has length \(1\).  Hence
\[
 \Gamma_1=\Delta_{\mathcal C_1}=\Gamma_0\ \cup\ h,
 \qquad h=[0,1]_\tau\ \text{attached to }\Gamma_0\text{ at }v_0,
\]
with \(\tau=v(w)\), \(\tau=0\) at \(v_0\) and \(\tau=1\) at \(x_E\).  The
curvature measure does not change: \(\rho_1^*\mathcal L\) has degree \(0\)
on \(E\) and degree \(1\) on each strict transform, so
\(\mu_0=\delta_{v_0}+\delta_{v_1}\) still, now viewed on \(\Gamma_1\); this
is \eqref{eq:ch21-curve-no-new-mass}.

On the hair, \(\lVert Z\rVert_0=|w|/\max\{1,|\pi w^2|,|w|\}=|w|\) for
\(|w|\leq1\), so
\begin{equation}
\label{eq:ch21-conic-restriction-1}
 \varphi|_{\Gamma_1}=
 \begin{cases}
  0, & \text{on }\Gamma_0,\\
  -\tau, & \text{on }h .
 \end{cases}
\end{equation}
Its curvature, by \eqref{eq:graph-laplacian}, is computed vertex by vertex:
at \(v_0\) the outgoing slopes are \(0\) into \(\Gamma_0\) and \(-1\) into
\(h\); at \(v_1\) the outgoing slope is \(0\); at \(x_E\) the outgoing slope
is \(+1\).  Hence
\begin{equation}
\label{eq:ch21-conic-curvature-1}
 \mu_0+\Delta_{\Gamma_1}\bigl(\varphi|_{\Gamma_1}\bigr)
 =(1-1)\delta_{v_0}+1\cdot\delta_{v_1}+1\cdot\delta_{x_E}
 =\delta_{v_1}+\delta_{x_E},
\end{equation}
of total mass \(2=\deg_CL\).  So one unit of mass has migrated from
\(v_0\) to the new vertex \(x_E\).  The corresponding model line bundle is
\(\rho_1^*\mathcal L-E\), whose degrees are
\((\rho_1^*\mathcal L-E)\cdot E=0-(-1)=1\),
\((\rho_1^*\mathcal L-E)\cdot\widetilde{C_0}=1-1=0\) and
\((\rho_1^*\mathcal L-E)\cdot\widetilde{C_1}=1-0=1\), confirming
\eqref{eq:ch21-conic-curvature-1} independently; and the corresponding
model function is \(\varphi_{-E}\), which takes the value \(-1\) at
\(x_E\), in accordance with \eqref{eq:ch20-sign}.

\emph{The tower.}  Repeating the blowup at the intersection of \(E\) with
the closure of \(\xi_0\), and simultaneously at \(\xi_1\) on the other side,
produces models \(\mathcal C_m\) with
\[
 \Gamma_m=\Gamma_0\cup h_0(m)\cup h_1(m),
\]
where \(h_j(m)\cong[0,m]\) is attached at \(v_j\) and carries the
coordinate \(\tau\).  On each hair \(\varphi=-\tau\), and
\begin{equation}
\label{eq:ch21-conic-curvature-m}
 \mu_0+\Delta_{\Gamma_m}\bigl(\varphi|_{\Gamma_m}\bigr)
 =\delta_{\tau=m\ \text{on}\ h_0}+\delta_{\tau=m\ \text{on}\ h_1},
\end{equation}
again of mass \(2\).  All the mass sits at the two extremities, at distance
\(m\) from \(\Gamma_0\).

\emph{Reading off the psh property.}  The retractions satisfy
\begin{equation}
\label{eq:ch21-conic-truncation}
 \varphi\circ p_{\mathcal C_m}=\max\{\varphi,-m\},
\end{equation}
because on a residue disc the retraction onto a segment of depth \(m\) sends
a point \(\zeta\) to the point of the segment of radius
\(\max\{|w(\zeta)|,|\pi|^m\}\).  The right-hand side of
\eqref{eq:ch21-conic-truncation} is a \(\theta\)-psh model function by
\eqref{eq:ch21-conic-curvature-m}, and these functions decrease to
\(\varphi\).  By \Cref{thm:ch21-decreasing-stability},
\(\varphi\in\PSH(C,\theta)\).  Alternatively, \(\varphi\) satisfies the two
clauses directly: it is upper semicontinuous as an infimum of continuous
functions, its restriction to \(\Gamma_m\) is exactly the model function
\eqref{eq:ch21-conic-restriction-1} extended, and clause~(1) is
\eqref{eq:ch21-conic-truncation}.

\emph{The polar set.}  From \eqref{eq:ch21-conic-truncation},
\(\varphi(\zeta)=-\infty\) exactly when \(|w(\zeta)|=0\), that is at
\(\xi_0\) and \(\xi_1\).  Hence
\[
 \{\varphi=-\infty\}=\{\xi_0,\xi_1\},
\]
two type-\(1\) points, in accordance with
\Cref{prop:ch21-polar-structure}\,\cref{itm:ch21-pol:curve}.  The measures
\eqref{eq:ch21-conic-curvature-m} converge weakly to
\(\delta_{\xi_0}+\delta_{\xi_1}\), of mass \(2\); this is a statement about
the model measures of the truncations only, and no measure is here attached
to \(\varphi\) itself, in accordance with \Cref{rem:ch21-polar-ma}.

\emph{What the example illustrates.}  The restriction of a psh function to
a dual complex may carry no information at all --- here
\eqref{eq:ch21-conic-restriction-0} is the zero function --- and become
informative only after the model is refined.  This is why
\Cref{def:ch21-bfj-psh} quantifies over all models and why
\Cref{prop:ch21-restriction-system}\,\cref{itm:ch21-rs:notsurjective} is not
a formality.
\end{example}

\subsection{A maximum and a convex combination, computed}
\label{sec:ch21-worked-max}

\begin{example}[Maximum and convex combination on the conic skeleton]
\label{ex:ch21-conic-numbers}
\index{maximum!of BFJ psh functions}
Keep \(C\), \(L\), \(\theta\), \(\Gamma_0=[0,1]\) and
\(\mu_0=\delta_{v_0}+\delta_{v_1}\) from
\Cref{ex:ch21-conic-refinement}.  A function \(u\) that is affine on
\(\Gamma_0\) with \(u(v_0)=A\), \(u(v_1)=B\) has
\[
 \mu_0+\Delta_{\Gamma_0}u=(1+B-A)\delta_{v_0}+(1+A-B)\delta_{v_1},
\]
which is \eqref{eq:running-perturbed-measure}; it is nonnegative exactly
when \(|A-B|\leq1\), which is \eqref{eq:running-nef-interval}.  Take the two
extreme cases
\[
 \varphi_1(s)=s-\tfrac12,\qquad \varphi_2(s)=\tfrac12-s
 \qquad(0\leq s\leq1),
\]
as in \Cref{ex:metrics-max-skeleton}, so that
\[
 \mu_0+\Delta\varphi_1=2\delta_{v_0},
 \qquad
 \mu_0+\Delta\varphi_2=2\delta_{v_1},
\]
both of mass \(2\).  Each lies on the boundary of the nef cone.  Both
\(\varphi_1\) and \(\varphi_2\) extend to \(\theta\)-psh model functions on
\(C^{\An}\), by \Cref{prop:ch21-curve-retraction} applied to their common
determining model, or directly by
\Cref{thm:ch21-model-consistency}.

\emph{The maximum.}  By
\Cref{thm:ch21-stability}\,\cref{itm:ch21-stab:max},
\(\varphi_{\max}:=\max\{\varphi_1,\varphi_2\}\in\PSH(C,\theta)\), and on
\(\Gamma_0\)
\[
 \varphi_{\max}(s)=\bigl|s-\tfrac12\bigr| .
\]
Its outgoing slopes are \(-1\) at \(v_0\), \(-1\) at \(v_1\), and the
slope jumps from \(-1\) to \(+1\) at \(s=\tfrac12\).  Hence
\begin{equation}
\label{eq:ch21-conic-max-measure}
 \mu_0+\Delta_{\Gamma_0}\varphi_{\max}
 =(1-1)\delta_{v_0}+2\,\delta_{1/2}+(1-1)\delta_{v_1}
 =2\,\delta_{1/2},
\end{equation}
of total mass \(2=\deg_CL\).  All the curvature has been swept to the
midpoint.  Concretely, \(2\delta_{1/2}\) is the Chambert--Loir measure of
the model obtained by blowing up the node once: the model
\(\{xy=\pi\}\) is replaced by a chain with a central component
\(E_{1/2}\) of multiplicity \(2\), and the corresponding line bundle has all
its relative degree on that component.

\emph{The convex combinations.}  By
\Cref{thm:ch21-stability}\,\cref{itm:ch21-stab:sum}, for \(t\in[0,1]\),
\[
 \varphi_t:=t\varphi_1+(1-t)\varphi_2
 =\bigl(t-\tfrac12\bigr)(2s-1)
\]
is again \(\theta\)-psh; it is affine with \(A=\varphi_t(v_0)=\tfrac12-t\)
and \(B=\varphi_t(v_1)=t-\tfrac12\), so \(B-A=2t-1\in[-1,1]\) and
\begin{equation}
\label{eq:ch21-conic-convex-measure}
 \mu_0+\Delta_{\Gamma_0}\varphi_t
 =2t\,\delta_{v_0}+(2-2t)\,\delta_{v_1},
\end{equation}
of total mass \(2\) for every \(t\).  The endpoints \(t=1\) and \(t=0\)
recover \(2\delta_{v_0}\) and \(2\delta_{v_1}\); the midpoint \(t=\tfrac12\)
gives \(\varphi_{1/2}\equiv0\) and returns the reference measure
\(\mu_0=\delta_{v_0}+\delta_{v_1}\).  Thus the convex combination
interpolates the curvature \emph{linearly}, while the maximum does something
else entirely.

\emph{The contrast, and a failure.}  The two operations of
\Cref{thm:ch21-stability} therefore act very differently on curvature:
\[
 \begin{array}{lll}
  \text{convex combination at }t=\tfrac12:
   &\varphi_{1/2}=0,
   &\mu=\delta_{v_0}+\delta_{v_1};\\[2pt]
  \text{maximum:}
   &\varphi_{\max}=|s-\tfrac12|,
   &\mu=2\delta_{1/2}.
 \end{array}
\]
Both have mass \(2\), and \(\varphi_{1/2}\leq\varphi_{\max}\) pointwise, as
it must be.  Finally, the \emph{minimum} is not psh, which is the failure
announced in \Cref{rem:ch21-not-stable}: for
\(\varphi_{\min}=\min\{\varphi_1,\varphi_2\}=-|s-\tfrac12|\) one finds
outgoing slopes \(+1\) at \(v_0\) and \(+1\) at \(v_1\) and a downward
kink of size \(-2\) at the midpoint, so that
\[
 \mu_0+\Delta_{\Gamma_0}\varphi_{\min}
 =2\delta_{v_0}-2\delta_{1/2}+2\delta_{v_1},
\]
of total mass \(2\) but with an atom of mass \(-2\).  The mass is right and
the sign is wrong; this is the standard way in which a plausible operation
fails, and it is why \Cref{prop:ch20-model-max} has a proof and
\(\min\) has none.
\end{example}

\subsection{A polar set computed}
\label{sec:ch21-worked-polar}

\Cref{ex:ch21-conic-refinement} already computed a polar set with two
points.  The next example computes an infinite one, on
\(\mathbf P^1\), and shows that a polar set need be neither closed nor
Zariski-closed.  It is the BFJ counterpart of
\Cref{ex:ch15-polar-nonanalytic}.

\begin{example}[A tower of models on \(\mathbf P^1\)]
\label{ex:ch21-p1-refinement}
Let \(X=\mathbf P^1_K\) with affine coordinate \(z\), \(L=\cO(1)\), and let
\(\|\cdot\|_0\) be the model metric of
\((\mathbf P^1_R,\cO(1))\), so that
\(\lVert s\rVert_0=|s|/\max\{|z_0|,|z_1|\}\) for a linear form \(s\).  Its
curvature form \(\theta\) has curvature measure
\(\delta_\eta\), of mass \(1=\deg\cO(1)\), where \(\eta=\zeta_{0,1}\) is the
Gauss point; the dual complex of \(\mathbf P^1_R\) is the single point
\(\{\eta\}\).

For \(a\in R\) let \(\mathcal X_m^{(a)}\) be the model obtained from
\(\mathbf P^1_R\) by \(m\) successive blowups along the closure of
\(\{z=a\}\), so that
\[
 \Delta_{\mathcal X^{(a)}_m}=[\eta,\ \zeta_{a,|\pi|^m}]
 \cong[0,m]_t,
 \qquad t=-\log r,
\]
the point of parameter \(t=j\) being \(\zeta_{a,|\pi|^{j}}\), a divisorial
point.  The curvature measure of \(\theta\) on this segment is
\(\delta_{t=0}\), of mass \(1\), because the pullback of \(\cO(1)\) has
degree \(0\) on every exceptional component.

Now take \(a=0\) and \(\varphi:=\log\lVert z_1\rVert_0\), where \(z_1\)
vanishes at \(z=0\).  In the affine chart,
\(\varphi=\log|z|-\log\max\{1,|z|\}\), so
\[
 \varphi=\log|z|\ \text{ on }\{|z|\leq1\},
 \qquad
 \varphi=0\ \text{ on }\{|z|\geq1\}.
\]
On the segment, \(\varphi(t)=-t\), and
\[
 \mu_0+\Delta_{[0,m]}\varphi
 =(1-1)\delta_{t=0}+1\cdot\delta_{t=m}
 =\delta_{\zeta_{0,|\pi|^m}},
\]
of mass \(1=\deg\cO(1)\).  Moreover
\(\varphi\circ p_{\mathcal X^{(0)}_m}=\max\{\varphi,-m\}\).  Hence
\(\varphi\in\PSH(\mathbf P^1,\theta)\) with polar set the single rational
point \(z=0\), and the model measures of its truncations converge weakly to
\(\delta_{\{z=0\}}\).

Two features are worth noting.  The tower \((\mathcal X^{(0)}_m)_m\) is
indexed by \(\mathbf N\), so on this tower sequences are legitimate; this is
the third mechanism of \Cref{rem:ch21-nets-scope}.  But the tower is
\emph{not} cofinal in the system of all models: the models
\(\mathcal X^{(a)}_1\) for \(a\) in distinct residue classes are not
dominated by any \(\mathcal X^{(0)}_m\), which is the content of
\Cref{lem:ch21-no-cofinal-sequence}.
\end{example}

\begin{example}[A polar set that is countably infinite and not closed]
\label{ex:ch21-polar-noncompact}
\index{polar set!BFJ}
Keep \(X=\mathbf P^1_K\), \(L=\cO(1)\), \(\|\cdot\|_0\) and \(\theta\) from
\Cref{ex:ch21-p1-refinement}, and assume the residue field
\(\widetilde K\) is infinite; for instance \(K=\mathbf C((t))\).  Choose
\(a_1,a_2,\ldots\in R\) lying in pairwise distinct residue classes, so that
\begin{equation}
\label{eq:ch21-separated-centres}
 |a_i-a_j|=1\qquad(i\neq j),
\end{equation}
and put \(c_j=2^{-j}\), so that \(\sum_{j\geq1}c_j=1\).  Let
\(\varphi_{a}:=\log\lVert z_1-az_0\rVert_0\), which by
\Cref{ex:ch21-p1-refinement} lies in \(\PSH(\mathbf P^1,\theta)\), is
\(\leq0\), and equals \(\log|z-a|\) on \(\{|z|\leq1\}\).  Define
\begin{equation}
\label{eq:ch21-polar-series}
 \Phi_N:=\sum_{j=1}^{N}c_j\varphi_{a_j}
        +\Bigl(\sum_{j>N}c_j\Bigr)\cdot 0,
 \qquad
 \varphi:=\inf_{N\geq1}\Phi_N .
\end{equation}

\emph{Each \(\Phi_N\) is \(\theta\)-psh.}  It is a convex combination of the
\(N+1\) functions \(\varphi_{a_1},\ldots,\varphi_{a_N},0\), each in
\(\PSH(\mathbf P^1,\theta)\) --- the last because \(\theta\) is semipositive
and \Cref{prop:ch21-constants}\,\cref{itm:ch21-const:const} applies --- so
\Cref{thm:ch21-stability}\,\cref{itm:ch21-stab:sum} gives
\(\Phi_N\in\PSH(\mathbf P^1,\theta)\).

\emph{The sequence decreases.}  \(\Phi_N-\Phi_{N+1}
=-c_{N+1}\varphi_{a_{N+1}}\geq0\), because
\(\varphi_{a}\leq0\).  Hence \(\varphi=\inf_N\Phi_N\) is the limit of a
decreasing sequence, itself a net.

\emph{The limit is not identically \(-\infty\).}  At the Gauss point
\(\eta\) one has \(|z-a_j|(\eta)=\max\{1,|a_j|\}=1\), so
\(\varphi_{a_j}(\eta)=0\) for every \(j\) and \(\varphi(\eta)=0\).  By
\Cref{thm:ch21-decreasing-stability},
\(\varphi\in\PSH(\mathbf P^1,\theta)\).

\emph{The polar set.}  Let \(b\in\mathbf P^1(K)\) be a rational point, that
is a point of type \(1\).  If \(|b|>1\) then \(\varphi(b)=0\).  If
\(|b|\leq1\), then by \eqref{eq:ch21-separated-centres} there is at most one
index \(j_0\) with \(|b-a_{j_0}|<1\), and \(|b-a_j|=1\) for all
\(j\neq j_0\).  Consequently
\[
 \varphi(b)=\sum_{j\geq1}c_j\log|b-a_j|
 =\begin{cases}
   c_{j_0}\log|b-a_{j_0}|, &\text{if such a }j_0\text{ exists},\\
   0, &\text{otherwise},
 \end{cases}
\]
a convergent expression in \([-\infty,0]\), and it equals \(-\infty\)
precisely when \(b=a_{j_0}\).  If instead \(x\) is a point of type \(2\),
\(3\) or \(4\), let \(r(x)>0\) be its radius in the classification
\Cref{thm:berk-a1-classification}; for \(x=\zeta_{b,r}\) one has
\(|z-a_j|(x)=\max\{|b-a_j|,r\}\geq r=r(x)\), and for a type-\(4\) point,
presented as a decreasing intersection of discs \(D(b_n,r_n)\) with
\(r_n\downarrow r(x)>0\), the same bound follows in the limit.  Hence
\(\varphi(x)\geq\bigl(\sum_jc_j\bigr)\log r(x)=\log r(x)>-\infty\) at every
such point.  Therefore
\begin{equation}
\label{eq:ch21-polar-computed}
 \{\varphi=-\infty\}=\{a_1,a_2,a_3,\ldots\},
\end{equation}
a countably infinite set of \(K\)-rational points.

\emph{It is not closed.}  A basic neighbourhood of the Gauss point \(\eta\)
is the complement of finitely many open discs, hence contains all but
finitely many \(a_j\).  So \(\eta\) lies in the closure of
\(\{a_j\}\), while \(\varphi(\eta)=0\).  \emph{It is not
Zariski-closed} either: an infinite subset of a curve that is
Zariski-closed is the whole curve.  Both conclusions are consistent with
\Cref{prop:ch21-polar-structure}: the polar set is a \(G_\delta\), it misses
every quasimonomial point, and it has empty interior, and nothing more is
claimed.

\emph{Two variants.}  Taking \(c_j=2^{-j}\) but allowing the \(a_j\) to
cluster --- for instance \(a_j\equiv a\bmod\pi\) for all \(j\) --- changes
the answer: the series \(\sum_jc_j\log|b-a_j|\) may then diverge at points
\(b\) that are not among the \(a_j\), and the polar set becomes larger and
harder to describe.  Taking instead finitely many terms returns a polar set
that is finite, hence closed and Zariski-closed.  The example is therefore
sharp in both directions: finiteness of the polar set is not automatic, and
its failure to be closed is produced by an entirely explicit series.
\end{example}

\subsection{Toric potentials}
\label{sec:ch21-worked-toric}

\begin{example}[Toric psh metrics]
\label{ex:ch21-toric}
\index{toric variety!BFJ psh metrics}
Let \(X=X_\Sigma\) be a polarized toric variety over \(K\), with lattice
\(N\), polytope \(P\subset M_{\mathbf R}\), and toric skeleton
\(N_{\mathbf R}\subset X^{\An}\) as in \Cref{sec:toric-convex}.  An
invariant model metric is
encoded by a rational piecewise affine convex function \(\phi\) on
\(N_{\mathbf R}\) with \(\phi-\check h_P\) bounded, the recession data being
fixed by the polarization.  Uniform completion gives continuous convex
functions with the same recession behavior, and decreasing completion gives
singular invariant psh metrics.  The maximum and convex-combination
operations in \Cref{thm:ch21-stability} become the ordinary
maximum and convex combination of convex functions.  This is an exact model
of the skeletal rules, but it does not replace the toroidal and multiplier-%
ideal arguments in the non-toric case.

Three points of the general theory become transparent here, and the reader
should check each against the corresponding general statement.

\emph{The retraction inequality is the statement that a convex function
attains its maximum on a face at a vertex.}  On a cone of \(\Sigma\) the
retraction to the toric skeleton is \(\trop\), the invariant potential is
constant on its fibres, and \eqref{eq:ch21-retraction-inequality} becomes
the standard inequality between a Laurent series and its tropicalization,
which is \Cref{lem:ch08-toric-gauss-section}.

\emph{Clause \textup{(2)} is uniform approximation of a convex function by
rational piecewise affine convex ones with the same recession function.}
This is the toric shadow of \Cref{def:ch21-bfj-psh}, and the approximation
is elementary: take the supremum of finitely many rational affine minorants
with slopes in \(-P\cap\tfrac1kM\), as in Step~3 of the proof of
\Cref{thm:ch19-toric-ma}.  The general case is not elementary, and it is
precisely the gap that \Cref{thm:ch21-model-consistency} fills.

\emph{Polar sets are visible.}  If \(\phi\) is allowed to take the value
\(-\infty\) on \(N_{\mathbf R}\), the resulting object is \emph{not} an
invariant BFJ psh function: by
\Cref{prop:ch21-first-consequences}\,\cref{itm:ch21-fc:finite} the
restriction to a dual complex is finite, and the toric skeleton is a union
of dual complexes.  Invariant psh functions therefore have polar sets
disjoint from \(N_{\mathbf R}\), and the polar set of an invariant psh
function lies in the ``boundary at infinity'' of the toric picture, exactly
as \(\{\xi_0,\xi_1\}\) did in \Cref{ex:ch21-conic-refinement}.

Finally, the toric case is where the difference between this chapter and
\Cref{part:cld} is smallest and therefore least instructive: for invariant
\emph{continuous} metrics the two theories agree by
\Cref{thm:ch19-toric-ma} and
\Cref{thm:ch19-cld-bfj-comparison}\,(B), and the singular invariant case is
not covered by either comparison.
\end{example}

\subsection{Changing the reference form}
\label{sec:ch21-worked-reference}

\begin{example}[Changing the reference form]
\label{ex:ch21-reference-class}
Suppose \(\theta\) is not semipositive but has ample generic class.  Choose
\(\psi\) as in \Cref{prop:ch20-positive-representative}.  Then

\begin{equation}
\label{eq:ch21-reference-class-example}
 \PSH(X,\theta)
   =\psi+\PSH(X,\theta+\ddc\psi).
\end{equation}

The right-hand reference form is positive on a model, so constants belong
to its psh class.  Formula \eqref{eq:ch21-reference-class-example} shows why
the theory naturally depends on an ample numerical class rather than on a
preselected semipositive representative.

The proof of \eqref{eq:ch21-reference-class-example} is the last paragraph
of the proof of \Cref{cor:ch21-reference-independence}: the map
\(u\mapsto u-\psi\) preserves upper semicontinuity, preserves clause~(1)
because \(\psi=\psi\circ p_{\mathcal X}\) for a determining model
(\Cref{lem:model-function-affine}), and carries
\(\mathcal P_{\mathcal X}(\theta)\) isometrically onto
\(\mathcal P_{\mathcal X}(\theta+\ddc\psi)\) by
\eqref{eq:ch20-reference-change}.

Two consequences are used repeatedly.  First, no generality is lost by
assuming \(\theta\) semipositive whenever a statement is invariant under
\eqref{eq:ch21-reference-class-example}; for instance the polar-set
statements of \Cref{prop:ch21-polar-structure} are, since \(\psi\) is
finite.  Second, generality \emph{is} lost for statements involving
normalizations by constants, since
\eqref{eq:ch21-reference-class-example} shifts a potential by the nonconstant
function \(\psi\); the normalization
\(\max_{\Delta_{\mathcal X}}\varphi=0\) of \Cref{chap:bfj-compactness} is
not preserved, and every energy statement in
\Cref{chap:bfj-energy} records its reference explicitly for this reason.
\end{example}

\section{Chapter dictionary}
\label{sec:ch21-dictionary}

\subsection{From the objects of this chapter to their use}
\label{sec:ch21-dict-forward}

\begin{center}
\small
\begin{tabular}{p{.27\textwidth}|p{.35\textwidth}|p{.30\textwidth}}
\toprule
\textbf{Object}&\textbf{Established property}&\textbf{First use}\\
\midrule
the class \(\PSH(X,\theta)\)
 &two clauses, cofinality
  (\Cref{def:ch21-bfj-psh},
   \Cref{prop:ch21-cofinal-independence})
 &\Cref{chap:bfj-compactness}\\
the skeletal closure
 \(\mathcal P_{\mathcal X}(\theta)\)
 &uniform closure of model psh restrictions
  \eqref{eq:ch21-skeletal-closure}
 &\Cref{thm:ch21-decreasing-stability},
  \Cref{thm:ch23-upper-envelope}\\
model consistency
 &\(\theta\)-psh model function \(\Leftrightarrow\) relatively nef
  (\Cref{thm:ch21-model-consistency})
 &\Cref{chap:bfj-ma},
  \Cref{thm:ch19-cld-bfj-comparison}\,(A)\\
closedness at divisorial points
 &\eqref{eq:ch21-model-closedness}
 &\Cref{thm:model-semipositive-closed}, and the proof of
  \Cref{thm:ch21-equivalent}\\
reconstruction
 &\(\varphi=\inf_{\mathcal X}\varphi\circ p_{\mathcal X}\)
  \eqref{eq:ch21-skeletal-reconstruction}
 &\Cref{chap:bfj-envelopes},
  \eqref{eq:ch23-upper-envelope-reconstruction}\\
divisorial determination
 &\Cref{thm:bfj-skeletal-characterization}\,(3),
  \Cref{prop:ch21-divisorial-order}
 &\Cref{chap:bfj-compactness},
  \Cref{chap:bfj-functoriality}\\
supremum on a skeleton
 &\eqref{eq:ch21-sup-on-skeleton}
 &normalizations in \Cref{chap:bfj-compactness}\\
stability calculus
 &constants, sums, maxima, uniform limits
  (\Cref{thm:ch21-stability})
 &everywhere in \Cref{part:bfj}\\
decreasing nets
 &\Cref{thm:ch21-decreasing-stability}
 &\Cref{thm:bfj-decreasing-regularization},
  \Cref{chap:bfj-ma}\\
pullback, model and continuous
 &\Cref{prop:ch21-pullback}
 &\Cref{chap:bfj-functoriality}\\
polar sets
 &\(G_\delta\), no quasimonomial point, empty interior
  (\Cref{prop:ch21-polar-structure})
 &\Cref{chap:bfj-capacity}\\
singularity types
 &partial order \(\preceq\)
  (\Cref{def:ch21-polar-singularity})
 &\Cref{chap:bfj-energy}\\
continuous comparison
 &Zhang semipositivity
  (\Cref{thm:ch21-continuous-zhang})
 &\Cref{thm:ch19-cld-bfj-comparison}\,(B),
  \Cref{chap:compare-ma}\\
\bottomrule
\end{tabular}
\end{center}

\subsection{From the later theory back to this chapter}
\label{sec:ch21-dict-backward}

\begin{center}
\small
\begin{tabular}{p{.30\textwidth}|p{.34\textwidth}|p{.28\textwidth}}
\toprule
\textbf{Later object}&\textbf{Smallest instance}&\textbf{Computed in}\\
\midrule
a singular \(\theta\)-psh function
 &\(\log\lVert Z\rVert_0\) on the conic; polar set two rational points
 &\Cref{ex:ch21-conic-refinement}\\
mass migrating to the polar set
 &\(\delta_{v_1}+\delta_{x_E}\to\delta_{\xi_0}+\delta_{\xi_1}\),
  mass \(2\)
 &\eqref{eq:ch21-conic-curvature-m}\\
maximum of two psh functions
 &\(|s-\tfrac12|\) on \([0,1]\), measure \(2\delta_{1/2}\)
 &\eqref{eq:ch21-conic-max-measure}\\
convex combination
 &\(2t\,\delta_{v_0}+(2-2t)\delta_{v_1}\)
 &\eqref{eq:ch21-conic-convex-measure}\\
failure of the minimum
 &an atom of mass \(-2\)
 &\Cref{ex:ch21-conic-numbers}\\
polar set not closed
 &\(\{a_1,a_2,\ldots\}\subset\mathbf P^1(K)\)
 &\eqref{eq:ch21-polar-computed}\\
no cofinal sequence of models
 &\(\mathbf P^1\) over \(\mathbf C((t))\)
 &\Cref{lem:ch21-no-cofinal-sequence}\\
clause~(1) is independent
 &\(\one_{\{x_0\}}\) at a rational point
 &\Cref{ex:ch21-clause-one-needed}\\
clause~(2) is independent
 &a concave function pulled back by \(p_{\mathcal X}\)
 &\Cref{ex:ch21-clause-two-needed}\\
\bottomrule
\end{tabular}
\end{center}

\subsection{Hypothesis ledger}
\label{sec:ch21-ledger}

\begin{center}
\small
\begin{tabular}{p{.30\textwidth}|p{.62\textwidth}}
\toprule
\textbf{Hypothesis or input}
 &\textbf{Where it is used, and what fails without it}\\
\midrule
\(\operatorname{char}\widetilde K=0\)
 &only in \Cref{thm:ch21-model-consistency}, through
  \eqref{eq:ch21-multiplier-subadditivity} and
  \eqref{eq:ch21-uniform-generation}; and in the cofinality of SNC models
  \Cref{prop:snc-cofinal}.  Without it the closedness statement
  \eqref{eq:ch21-model-closedness} is not available by this proof\\
\(K\) discretely valued
 &everywhere: the index set of \Cref{def:ch21-bfj-psh} is the set of models
  over \(R\), and \(\varphi_{\mathcal X_s}=1\) normalizes divisorial points\\
\(\{\theta\}\) ample
 &only in \Cref{prop:ch21-constants}\,\cref{itm:ch21-const:nonempty}, for
  nonemptiness; every other statement is conditional on membership\\
\(\theta\) semipositive
 &never assumed; used as a hypothesis in
  \Cref{prop:ch21-constants}\,\cref{itm:ch21-const:const} and removable by
  \eqref{eq:ch21-reference-class-example}\\
\Cref{thm:bfj-skeletal-equicontinuity}
 &\Cref{thm:ch21-decreasing-stability} and the implication
  \cref{itm:ch21-eq:divisorial}\(\Rightarrow\)\cref{itm:ch21-eq:inf} of
  \Cref{thm:ch21-equivalent}; without it a pointwise divisorial limit gives
  no control on a whole dual complex\\
\Cref{thm:bfj-decreasing-regularization}
 &only in \Cref{thm:ch21-continuous-zhang} and
  \Cref{prop:ch21-pullback}\,\cref{itm:ch21-pb:cont}; it is a theorem of
  \Cref{chap:bfj-compactness} and is never used to prove anything it
  depends on\\
\(\dim X=1\)
 &\Cref{prop:ch21-curve-retraction} and
  \Cref{prop:ch21-polar-structure}\,\cref{itm:ch21-pol:curve} only; both are
  false or unproved in higher dimension\\
\bottomrule
\end{tabular}
\end{center}

\subsection{Dependency chain}
\label{sec:ch21-chain}

\[
 \begin{aligned}
  &\text{skeletal calculus of \Cref{chap:skeleta}}
   \Longrightarrow
   \Cref{lem:ch21-model-skeletal-properties};\\
  &\Cref{lem:ch21-model-skeletal-properties}
   +\ \text{multiplier ideals}
   \Longrightarrow
   \Cref{thm:ch21-model-consistency};\\
  &\Cref{thm:ch21-model-consistency}
   \Longrightarrow
   \Cref{cor:ch21-reference-independence},\
   \Cref{prop:ch21-constants};\\
  &\text{clause (1)}+\text{usc}
   \Longrightarrow
   \Cref{thm:bfj-skeletal-characterization};\\
  &\Cref{thm:bfj-skeletal-equicontinuity}
   \Longrightarrow
   \Cref{thm:ch21-decreasing-stability},\
   \Cref{thm:ch21-equivalent};\\
  &\Cref{thm:bfj-decreasing-regularization}
   \Longrightarrow
   \Cref{thm:ch21-continuous-zhang}
   \Longrightarrow
   \Cref{prop:ch21-pullback}\,\cref{itm:ch21-pb:cont}.
 \end{aligned}
\]

Only the second, fifth and sixth lines use theorems proved elsewhere in the
book --- they are the three imports of \Cref{sec:ch21-imported} --- and the
first uses the skeletal calculus of \Cref{chap:skeleta}; the remaining lines
are self-contained in this chapter.

\section*{Notes and further reading}

\paragraph{The definition and its immediate theory.}
The definition, model-consistency theorem, skeletal reconstruction,
stability statements, and topology originate in
\cite[Sections~5--7]{BFJ16}.  Our \Cref{def:ch21-bfj-psh} is their
definition verbatim; the reorganization here consists in separating the two
clauses logically (\Cref{ex:ch21-clause-one-needed},
\Cref{ex:ch21-clause-two-needed}), in isolating the skeletal closure
\eqref{eq:ch21-skeletal-closure} as a named object so that clause~(2) can be
read as a membership, and in collecting the equivalent formulations into a
single statement, \Cref{thm:ch21-equivalent}.  None of these is a new
theorem.  The precursor of the whole circle of ideas is the earlier
valuative work of the same authors on singularities of plurisubharmonic
functions \cite{BFJ08,BFJ09}, where the passage from divisorial data to
functions on a space of valuations is first carried out systematically.

\paragraph{The closedness theorem.}
The closedness theorem used in
\Cref{thm:ch21-model-consistency} is one of the places where the original
residue-characteristic-zero proof uses multiplier ideals; the package we
import is \cite[Appendix~B]{BFJ16}, and the underlying asymptotic
multiplier-ideal theory is Lazarsfeld's \cite[Chapter~11]{PAG2}, with the
subadditivity theorem going back to Demailly, Ein and Lazarsfeld
\cite{DEL00}.  The statement that a model metric which is a pointwise limit
at divisorial points of semipositive model metrics is itself semipositive is
proved independently, over an arbitrary complete nontrivially valued field,
as \Cref{thm:model-semipositive-closed}; the reader comparing the two proofs
will see that the multiplier-ideal input is what allows the present chapter
to reach the same conclusion while simultaneously producing the
approximating net from clause~(2).

\paragraph{Compactness, regularization, and envelopes.}
Compactness and
regularization are separated into \Cref{chap:bfj-compactness} and
\Cref{chap:bfj-envelopes}, following the logical order of
\cite[Sections~6--8]{BFJ16} rather than its order of exposition.  The
equicontinuity estimate \Cref{thm:bfj-skeletal-equicontinuity} is
\cite[Theorem~6.1]{BFJ16}; the decreasing regularization theorem
\Cref{thm:bfj-decreasing-regularization} is \cite[Theorem~8.7]{BFJ16}.  We
have deliberately not used regularization to prove anything in the present
chapter except \Cref{thm:ch21-continuous-zhang} and the continuous case of
pullback, so that the reader can see that the class \(\PSH(X,\theta)\) has a
self-contained elementary theory before the hard approximation theorems
arrive.

\paragraph{The comparison with Zhang's metrics.}
The continuous comparison with Zhang's metrics
is \cite[Corollary~8.8]{BFJ16}; Zhang's original definition of a
semipositive continuous metric is \cite{Zhang1995}, and the systematic
comparison of the algebraic, formal and piecewise linear notions of
semipositivity, over general complete fields, is Gubler and Martin
\cite{GM19}.  \Cref{thm:gm-metric-comparison} records the part of that
comparison used in this book.  Pullback beyond the model locus is treated,
with its hypotheses made explicit, in \Cref{chap:bfj-functoriality}; it is
not silently included among the elementary closure operations here.

\paragraph{The distinction from the CLD class.}
The insistence throughout this chapter on separating global approximation
from local regularity is not a stylistic preference.  The local theory of
Chambert-Loir and Ducros \cite{CLD25} defines psh functions as sections of a
sheaf generated by local decreasing limits, and its first stage is the class
of \emph{locally approachable} functions,
\cite[Definition~8.1.13]{CLD25}, restated here as
\Cref{def:ch15-approachable}.  Whether that hierarchy collapses at the first
stage is open in general, and it is the reason no formal argument passes from
BFJ to CLD in higher dimension; see
\Cref{thm:ch15-approachable-stability} and item (L5) of
\Cref{sec:ch19-limitations-table}.  In the separate scope where $K$ is
algebraically closed, complete, and nontrivially valued and $C$ is a smooth
proper algebraic curve, Wanner \cite{Wanner19} identifies current positivity
on $C^{\An}$ with Thuillier subharmonicity, as reproduced in
\Cref{thm:wanner-curve-comparison}, and
\Cref{thm:p063-current-psh-ordinary-regularization} identifies that
continuous current-positive class locally with the first ordinary-smooth
decreasing-limit stage.  This supplies neither descent to the standing BFJ
base field nor a formal identification of an arbitrary global BFJ class with
the CLD class.  The metric comparison proved in this book is
\Cref{thm:ch19-cld-bfj-comparison}, and
\Cref{sec:ch21-cld-contrast} states its scope in the form used by
\Cref{part:bfj}.  A reader who wants the two classes compared through
positivity of currents rather than through metrics should consult
\cite{GK17,GK19} for the delta-form formulation and
\cite{GR25Classical,GR25PL} for the curve-test formulation.

\paragraph{Nets.}
\Cref{lem:ch21-no-cofinal-sequence} is elementary and is included because
the assertion that the model system admits no countable cofinal family is
often stated without proof.  The proof given is the one-dimensional case;
the higher-dimensional statement follows by pulling back along a pencil, but
we have not needed it, since a counterexample on a curve already blocks any
general sequential reduction.  The abstract counterparts are
\Cref{ex:no-cofinal-sequence} and
\Cref{ex:decreasing-net-not-sequential}, both in
\Cref{chap:topological-tools}, and the checklist governing every sequential
reduction in this book is \Cref{rem:net-vs-sequence-checklist}.

\paragraph{What is new here, and what is not.}
No theorem of this chapter is new.  The statements that do not appear in
this form in \cite{BFJ16} are: the independence examples
\Cref{ex:ch21-clause-one-needed} and
\Cref{ex:ch21-clause-two-needed}; the packaging
\Cref{thm:ch21-equivalent}; the curve retraction statement
\Cref{prop:ch21-curve-retraction}, which is elementary but is used in
\Cref{sec:ch21-worked} and delimits what fails in higher dimension; the
polar-set structure \Cref{prop:ch21-polar-structure} and the explicit
example \Cref{ex:ch21-polar-noncompact}; the separations
\Cref{prop:ch21-bfj-vs-cld}; and \Cref{lem:ch21-no-cofinal-sequence}.  All
are either elementary consequences of the cited results or explicit
computations, and no priority is claimed for any of them.  The worked
examples of \Cref{sec:ch21-worked} are carried to numbers on purpose: every
mass appearing there is checked against \(\deg_CL\), and the same
computations are reused in \Cref{chap:bfj-examples}.
